%%%%%%%%%%%%%%%%%%%%%%%%%%%%%%%%%%%%%%%%%%%%%%%%%%%%%%%%%%%%%%%%%%%%%%%%%%%%%%%%%%
\begin{frame}[fragile]\frametitle{}
\begin{center}
{\Large Yoga for Health Promotion }
\end{center}
\end{frame}


%%%%%%%%%%%%%%%%%%%%%%%%%%%%%%%%%%%%%%%%%%%%%%%%%%%%%%%%%%%
\begin{frame}[fragile]\frametitle{Syllabus}

\begin{itemize}
\item 3.1  Brief introduction to human body. 
\item 3.2  Meaning and Means of health promotion and role of Yoga  in health promotion. 
\item 3.3  Yogic positive attitudes ( Maîtri, Karuna, Mudita, Upeksha). 
\item 3.4  Concept of bhavas (Dharma, Jnana, Vairagya, Aishvarya) and their relevance in well being. 
\item 3.5  Dincharya and Ritucharya with respect to Yogic life style. 
\item 3.6  Holistic approach of Yoga towards health and diseases. 
\item 3.7  Introduction to First aid and Cardio Pulmonary Resuscitation (CPR). 
\item 3.8  Yogic management of stress and its consequences. 
\item 3.9  Yoga in prevention of metabolic and respiratory disorders. 
\item 3.10  Yoga for personality development. 
\end{itemize}
	  
\end{frame}

%%%%%%%%%%%%%%%%%%%%%%%%%%%%%%%%%%%%%%%%%%%%%%%%%%%%%%%%%%%%%%%%%%%%%%%%%%%%%%%%%%
\begin{frame}[fragile]\frametitle{}
\begin{center}
{\Large 3.1  Brief introduction to human body: Anatomy}
\end{center}
\end{frame}

%%%%%%%%%%%%%%%%%%%%%%%%%%%%%%%%%%%%%%%%%%%%%%%%%%%%%%%%%%%
\begin{frame}[fragile]\frametitle{Body System, Organs and their Functions}

		\begin{center}
		% FIG ycb1-fig-bodysystems | now: images/ycb1_bodysystems.png | final: no image file (see register) | register: books/YCB1/docs/image_register.csv
		% PROMPT SP-NONE (common text: books/YCB1/docs/image_system_prompt.md) | Do not generate a picture. This is a screenshot of a table. Retype it as a LaTeX table (deck) and a Markdown table (notes) from the checked notes block ycb1-hl-human-body-introduction, in own words where it is a list of facts. Used in block ycb1-hl-human-body-introduction, frame 'Body System, Organs and their Functions'.
		% Screenshot retired; retyped from the notes Material bullets.
		% \includegraphics[width=\linewidth,keepaspectratio]{ycb1_bodysystems}
		{\scriptsize
		\begin{tabular}{@{}p{0.16\linewidth}p{0.38\linewidth}p{0.4\linewidth}@{}}
		\toprule
		\textbf{System} & \textbf{Main parts} & \textbf{Main work} \\
		\midrule
		Skeletal & Bones, cartilage, ligaments, tendons, joints & Framework, support, protection; movement; blood cells in the marrow \\
		Muscular & Voluntary muscles; involuntary muscles of heart and gut & Moves the body and the organs \\
		Nervous & Brain, spinal cord, nerves; sense organs & Control and quick messages; autonomic part works without our will \\
		Endocrine & Glands & Hormones that control growth, mood and energy \\
		Circulatory & Heart, blood vessels & Carries blood, oxygen, nutrients and waste \\
		Respiratory & Nose, windpipe, lungs; diaphragm & Takes in oxygen, gives out carbon dioxide \\
		Digestive & The gut & Breaks down food, takes in nutrients \\
		Excretory & Kidneys, bladder; skin, lungs & Removes waste; balances water and salts \\
		Reproductive & Reproductive organs & Continuation of the species \\
		\bottomrule
		\end{tabular}}
		\end{center}

\end{frame}

%%%%%%%%%%%%%%%%%%%%%%%%%%%%%%%%%%%%%%%%%%%%%%%%%%%%%%%%%%%
\begin{frame}[fragile]\frametitle{Components of the Skeletal System}
\begin{columns}
    \begin{column}[T]{0.5\linewidth}
      \begin{itemize}
		\item Bones: about 206 in an adult; living tissue with a good blood supply
		\item Cartilage: elastic tissue that cushions the joints
		\item Ligaments: join bone to bone
		\item Tendons: join muscle to bone
		\item Joints: places where bones meet
	  \end{itemize}
    \end{column}
    \begin{column}[T]{0.5\linewidth}
		\begin{center}
		% FIG ycb1-fig-skeleton | now: images/skeleton.jpg | final: books/YCB1/images/new/skeleton_new.pdf | register: books/YCB1/docs/image_register.csv
		% PROMPT SP-DIAGRAM (common text: books/YCB1/docs/image_system_prompt.md) | Used in block ycb1-hl-human-body-introduction, frame 'Components of the Skeletal System'. Anatomy diagram: Components of the Skeletal System. A simple schematic in front or side view, whichever shows the structure best; main parts as flat shapes with thin outlines, grayscale-safe. Labels: only the terms listed under this figure in the notes block (list them here before generating); each label short, horizontal, with a leader line. One idea only; leave out finer detail. Aspect ratio 4:3. Check before use: qualified subject reviewer (health and anatomy facts).
		% NEW-PIC (uncomment to show the new picture beside the original; comment the original only after review) \includegraphics[width=0.35\linewidth,keepaspectratio]{new/skeleton_new.pdf}
		\includegraphics[width=\linewidth,keepaspectratio]{skeleton.jpg} 
		
		{\tiny (Ref: https://www.swiss-alp-health.ch/en/what-is-a-joint/)}
		\end{center}	

    \end{column}
  \end{columns}
\end{frame}



%%%%%%%%%%%%%%%%%%%%%%%%%%%%%%%%%%%%%%%%%%%%%%%%%%%%%%%%%%%
\begin{frame}[fragile]\frametitle{Divisions of the Skeletal System}
\begin{columns}
    \begin{column}[T]{0.5\linewidth}
      \begin{itemize}
		\item Axial skeleton: skull, vertebral column, rib cage
			\begin{itemize}
			\item Skull: 22 bones; guards the brain, eyes and inner ear
			\item Spine: a column of vertebrae; holds up the trunk, guards the spinal cord
			\item Rib cage: 12 pairs of ribs; guards the lungs and heart
			\end{itemize}

		\item Appendicular skeleton: shoulder and pelvic girdles, and the limbs
			\begin{itemize}
			\item Shoulder girdle: shoulder blades and collar bones
			\item Arm: humerus, radius, ulna, wrist bones, hand bones, finger bones
			\item Pelvis: flat bones joined to the sacrum; the base for the legs
			\item Leg: femur, patella, tibia, fibula, ankle bones, foot bones, toe bones
		    \end{itemize}
	  \end{itemize}
    \end{column}
    \begin{column}[T]{0.5\linewidth}
		\begin{center}
		% FIG ycb1-fig-skeleton-divisions | now: images/skeleton_divisions.jpg | final: books/YCB1/images/new/skeleton_divisions_new.pdf | register: books/YCB1/docs/image_register.csv
		% PROMPT SP-DIAGRAM (common text: books/YCB1/docs/image_system_prompt.md) | Used in block ycb1-hl-human-body-introduction, frame 'Divisions of the Skeletal System'. Anatomy diagram: Divisions of the Skeletal System. A simple schematic in front or side view, whichever shows the structure best; main parts as flat shapes with thin outlines, grayscale-safe. Labels: only the terms listed under this figure in the notes block (list them here before generating); each label short, horizontal, with a leader line. One idea only; leave out finer detail. Aspect ratio 4:3. Check before use: qualified subject reviewer (health and anatomy facts).
		% NEW-PIC approved by the author (2026-11-27); the original is commented out below it
		\includegraphics[width=\linewidth,keepaspectratio]{new/skeleton_divisions_new.pdf}
		% original, replaced (author approved 2026-11-27): \includegraphics[width=\linewidth,keepaspectratio]{skeleton_divisions.jpg}
		\end{center}	
    \end{column}
  \end{columns}
\end{frame}

%%%%%%%%%%%%%%%%%%%%%%%%%%%%%%%%%%%%%%%%%%%%%%%%%%%%%%%%%%%
\begin{frame}[fragile]\frametitle{Functions of skeleton system}

\begin{itemize}
\item Structural Framework
\item Support and protection
\item Making blood cells, in the marrow
\item Storehouse of minerals
\end{itemize}
	  
\end{frame}


%%%%%%%%%%%%%%%%%%%%%%%%%%%%%%%%%%%%%%%%%%%%%%%%%%%%%%%%%%%
\begin{frame}[fragile]\frametitle{Vertebral Column}
\begin{columns}
    \begin{column}[T]{0.65\linewidth}
      \begin{itemize}
		\item Spine: a column of vertebrae from the skull to the lower back
		\item 33 vertebrae, joined by cartilage and ligaments
		\item The spinal cord runs in the central canal
		\item Regions from the top: cervical (7), thoracic (12), lumbar (5)
		\item Then sacral (5 fused into 1) and coccygeal (4 fused into 1)
		\item Natural curves: cervical, thoracic, lumbar, sacral (pelvic)
		\item Faulty posture can exaggerate the curves
		\item Kyphosis: too much thoracic curve
		\item Lordosis: too much lumbar curve
		\item Scoliosis: a sideways curve of the spine
		\item Asanas strengthen and stretch muscles and move the spine, which supports good posture
	  \end{itemize}
    \end{column}
    \begin{column}[T]{0.35\linewidth}
		\begin{center}
		% FIG ycb1-fig-vertebral-column | now: images/vertebral_column.jpg | final: books/YCB1/images/new/vertebral_column_new.pdf | register: books/YCB1/docs/image_register.csv
		% PROMPT SP-DIAGRAM (common text: books/YCB1/docs/image_system_prompt.md) | Used in block ycb1-hl-human-body-introduction, frame 'Vertebral Column'. Anatomy diagram: Vertebral Column. A simple schematic in front or side view, whichever shows the structure best; main parts as flat shapes with thin outlines, grayscale-safe. Labels: only the terms listed under this figure in the notes block (list them here before generating); each label short, horizontal, with a leader line. One idea only; leave out finer detail. Aspect ratio 4:3. Check before use: qualified subject reviewer (health and anatomy facts).
		% NEW-PIC approved by the author (2026-11-27); the original is commented out below it
		\includegraphics[width=\linewidth,keepaspectratio]{new/vertebral_column_new.pdf}
		% original, replaced (author approved 2026-11-27): \includegraphics[width=\linewidth,keepaspectratio]{vertebral_column.jpg}
		\end{center}	
    \end{column}
  \end{columns}
\end{frame}

%%%%%%%%%%%%%%%%%%%%%%%%%%%%%%%%%%%%%%%%%%%%%%%%%%%%%%%%%%%
\begin{frame}[fragile,shrink=6]\frametitle{Types of Spinal Movements}
\begin{columns}
    \begin{column}[T]{0.5\linewidth}
      \begin{itemize}
		\item Flexion: forward bend; largest in the neck (e.g., Paschimottanasana, Padahastasana)
		\item Extension: back bend (e.g., Bhujangasana, Dhanurasana)
		\item Rotation: twist; largest between the first two neck vertebrae, atlas and axis (e.g., Ardha Matsyendrasana)
		\item Side bending: largest in the neck and lower back (e.g., Trikonasana)
		% Circumduction bullet removed: the notes' checked list has five movements only, and the spelling "Chakki Chalavan" is unsettled.
		% \item Circumduction: Swaying combining all movements (e.g., Chakki Chalavan; an addition to the guidebook's five movements)
		\item Elongation: stretch upwards from the base of the spine (e.g., Tadasana, Urdhva Hastasana)
	  \end{itemize}
    \end{column}
    \begin{column}[T]{0.5\linewidth}
		\begin{center}
		% FIG ycb1-fig-spinal-movements | now: images/spinal_movements.jpg | final: books/YCB1/images/new/spinal_movements_new.pdf | register: books/YCB1/docs/image_register.csv
		% PROMPT SP-DIAGRAM (common text: books/YCB1/docs/image_system_prompt.md) | Used in block ycb1-hl-human-body-introduction, frame 'Types of Spinal Movements'. Anatomy diagram: Types of Spinal Movements. A simple schematic in front or side view, whichever shows the structure best; main parts as flat shapes with thin outlines, grayscale-safe. Labels: only the terms listed under this figure in the notes block (list them here before generating); each label short, horizontal, with a leader line. One idea only; leave out finer detail. Aspect ratio 4:3. Check before use: qualified subject reviewer (health and anatomy facts).
		% NEW-PIC (uncomment to show the new picture beside the original; comment the original only after review) \includegraphics[width=0.35\linewidth,keepaspectratio]{new/spinal_movements_new.pdf}
		\includegraphics[width=\linewidth,keepaspectratio]{spinal_movements.jpg} 
		\end{center}	
    \end{column}
  \end{columns}
\end{frame}

%%%%%%%%%%%%%%%%%%%%%%%%%%%%%%%%%%%%%%%%%%%%%%%%%%%%%%%%%%%
\begin{frame}[fragile]\frametitle{Types of Joints}

      \begin{itemize}
		\item Joint: a place where two bones meet
		\item Fibrous: little or no movement; e.g., the sutures of the skull
		\item Cartilaginous: bones joined by cartilage; strong and slightly flexible
		\item Example: ribs to breastbone, which lets the chest move in breathing
		\item Synovial: the most mobile; bone ends coated with cartilage
		\item Synovial ends sit in a sealed capsule filled with fluid
		\item Synovial joints come in several kinds, each for certain movements
	  \end{itemize}

\end{frame}

%%%%%%%%%%%%%%%%%%%%%%%%%%%%%%%%%%%%%%%%%%%%%%%%%%%%%%%%%%%
\begin{frame}[fragile]\frametitle{Types of Joints}

		\begin{center}
		% FIG ycb1-fig-joints | now: images/joints.jpg | final: books/YCB1/images/new/joints_new.pdf | register: books/YCB1/docs/image_register.csv
		% PROMPT SP-DIAGRAM (common text: books/YCB1/docs/image_system_prompt.md) | Used in block ycb1-hl-human-body-introduction, frame 'Types of Joints'. Anatomy diagram: Types of Joints. A simple schematic in front or side view, whichever shows the structure best; main parts as flat shapes with thin outlines, grayscale-safe. Labels: only the terms listed under this figure in the notes block (list them here before generating); each label short, horizontal, with a leader line. One idea only; leave out finer detail. Aspect ratio 4:3. Check before use: qualified subject reviewer (health and anatomy facts).
		% NEW-PIC (uncomment to show the new picture beside the original; comment the original only after review) \includegraphics[width=0.35\linewidth,keepaspectratio]{new/joints_new.pdf}
		\includegraphics[width=0.5\linewidth,keepaspectratio]{joints.jpg} 
		\end{center}	

\end{frame}

%%%%%%%%%%%%%%%%%%%%%%%%%%%%%%%%%%%%%%%%%%%%%%%%%%%%%%%%%%%
\begin{frame}[fragile]\frametitle{Muscular System Overview}

      \begin{itemize}
		\item Muscles are tissues that contract.
		\item They turn chemical energy into movement.
		\item Three types: voluntary, involuntary and cardiac.
		\item Voluntary: under conscious control; used in walking, balancing, writing.
		\item Involuntary: run by the autonomic nervous system; used in digestion and blood flow.
		\item Cardiac: found only in the heart; keeps its own rhythm.
	  \end{itemize}

\end{frame}

%%%%%%%%%%%%%%%%%%%%%%%%%%%%%%%%%%%%%%%%%%%%%%%%%%%%%%%%%%%
\begin{frame}[fragile]\frametitle{Muscular System Overview}

		\begin{center}
		% FIG ycb1-fig-muscle-chart-showing-the-muscular-system-labeled | now: images/Muscle-chart-showing-the-muscular-system-labeled.jpg | final: books/YCB1/images/new/Muscle-chart-showing-the-muscular-system-labeled_new.pdf | register: books/YCB1/docs/image_register.csv
		% PROMPT SP-DIAGRAM (common text: books/YCB1/docs/image_system_prompt.md) | Used in block ycb1-hl-human-body-introduction, frame 'Muscular System Overview'. Anatomy diagram: Muscular System Overview. A simple schematic in front or side view, whichever shows the structure best; main parts as flat shapes with thin outlines, grayscale-safe. Labels: only the terms listed under this figure in the notes block (list them here before generating); each label short, horizontal, with a leader line. One idea only; leave out finer detail. Aspect ratio 4:3. Check before use: qualified subject reviewer (health and anatomy facts).
		% NEW-PIC (uncomment to show the new picture beside the original; comment the original only after review) \includegraphics[width=0.35\linewidth,keepaspectratio]{new/Muscle-chart-showing-the-muscular-system-labeled_new.pdf}
		\includegraphics[width=0.6\linewidth,keepaspectratio]{Muscle-chart-showing-the-muscular-system-labeled.jpg}
						
		{\tiny (Ref: https://www.biologyonline.com/dictionary/muscular-system)}			
		\end{center}	

\end{frame}


%%%%%%%%%%%%%%%%%%%%%%%%%%%%%%%%%%%%%%%%%%%%%%%%%%%%%%%%%%%
\begin{frame}[fragile]\frametitle{Functions of Muscular System}

      \begin{itemize}
		\item Movement of the body
		\item Posture: holding the body up against gravity
		\item Protection of the internal organs
		\item Heat production
		% Removed: not in the notes' checked list of muscle functions (stability of joints replaces it).
		% \item Store for energy (protein and carbohydrates)
		\item Stability of the joints
		\item Movement inside organs, by smooth and heart muscle
	  \end{itemize}

\end{frame}

%%%%%%%%%%%%%%%%%%%%%%%%%%%%%%%%%%%%%%%%%%%%%%%%%%%%%%%%%%%
\begin{frame}[fragile]\frametitle{Muscle Contraction}

      \begin{itemize}
		\item A muscle is made of fibres wrapped in a sheath.
		\item Each fibre holds thin (actin) and thick (myosin) filaments.
		\item In contraction the filaments slide over each other and the fibre shortens.
		\item At rest the filaments overlap little.
		\item When stimulated they slide further over each other; at full contraction the overlap is greatest.
		\item Strength grows when more fibres take part, or work better.
		\item Isotonic: the muscle changes length while the load stays the same.
		\item Isometric: the muscle keeps its length while the tension rises.
	  \end{itemize}

\end{frame}

%%%%%%%%%%%%%%%%%%%%%%%%%%%%%%%%%%%%%%%%%%%%%%%%%%%%%%%%%%%
\begin{frame}[fragile]\frametitle{Muscle Contraction}

		\begin{center}
		% FIG ycb1-fig-mechnaism-of-muscle-contrtaction | now: images/Mechnaism-of-muscle-contrtaction.png | final: books/YCB1/images/new/Mechnaism-of-muscle-contrtaction_new.pdf | register: books/YCB1/docs/image_register.csv
		% PROMPT SP-DIAGRAM (common text: books/YCB1/docs/image_system_prompt.md) | Used in block ycb1-hl-human-body-introduction, frame 'Muscle Contraction'. Anatomy diagram: Muscle Contraction. A simple schematic in front or side view, whichever shows the structure best; main parts as flat shapes with thin outlines, grayscale-safe. Labels: only the terms listed under this figure in the notes block (list them here before generating); each label short, horizontal, with a leader line. One idea only; leave out finer detail. Aspect ratio 4:3. Check before use: qualified subject reviewer (health and anatomy facts).
		% NEW-PIC (uncomment to show the new picture beside the original; comment the original only after review) \includegraphics[width=0.35\linewidth,keepaspectratio]{new/Mechnaism-of-muscle-contrtaction_new.pdf}
		\includegraphics[width=0.8\linewidth,keepaspectratio]{Mechnaism-of-muscle-contrtaction.png}
		
						
		{\tiny (Ref: https://www.geeksforgeeks.org/mechanism-of-muscle-contraction-class-11/)}		
		\end{center}	

\end{frame}

%%%%%%%%%%%%%%%%%%%%%%%%%%%%%%%%%%%%%%%%%%%%%%%%%%%%%%%%%%%
\begin{frame}[fragile]\frametitle{Reflex Action \& Reciprocal Inhibition}

      \begin{itemize}
		% Removed: not in the notes, and a motor unit is one nerve cell with the fibres it supplies, not a fibre group.
		% \item Motor units: smallest nerve fibre groups in muscles.
		\item Proprioceptors: sensors that tell the brain where the body is.
		\item This sense helps posture and coordination.
		\item Stretch reflex: a quick contraction when a muscle is suddenly lengthened.
		\item Example: back muscles tighten when one bends forward quickly.
		\item Slow movements support deep breathing.
		\item Reciprocal inhibition: when one muscle contracts, its opposite relaxes.
		\item Example: the biceps contract, the triceps relax.
		\item This gives smooth, coordinated movement.
	  \end{itemize}

\end{frame}



%%%%%%%%%%%%%%%%%%%%%%%%%%%%%%%%%%%%%%%%%%%%%%%%%%%%%%%%%%%
\begin{frame}[fragile]\frametitle{Reflex Action \& Reciprocal Inhibition}

		\begin{center}
		% FIG ycb1-fig-reciprocal-inhibition | now: images/Reciprocal-Inhibition.jpg | final: books/YCB1/images/new/Reciprocal-Inhibition_new.pdf | register: books/YCB1/docs/image_register.csv
		% PROMPT SP-DIAGRAM (common text: books/YCB1/docs/image_system_prompt.md) | Used in block ycb1-hl-human-body-introduction, frame 'Reflex Action and Reciprocal Inhibition'. Anatomy diagram: Reflex Action and Reciprocal Inhibition. A simple schematic in front or side view, whichever shows the structure best; main parts as flat shapes with thin outlines, grayscale-safe. Labels: only the terms listed under this figure in the notes block (list them here before generating); each label short, horizontal, with a leader line. One idea only; leave out finer detail. Aspect ratio 4:3. Check before use: qualified subject reviewer (health and anatomy facts).
		% NEW-PIC approved by the author (2026-11-27); the original is commented out below it
		\includegraphics[width=0.8\linewidth,keepaspectratio]{new/Reciprocal-Inhibition_new.pdf}
		% original, replaced (author approved 2026-11-27): \includegraphics[width=0.8\linewidth,keepaspectratio]{Reciprocal-Inhibition}
		
						
		% credit of the replaced original: {\tiny (Ref: https://www.corewalking.com/reciprocal-inhibition/)}
		{\tiny (Line diagram drawn for this deck.)}
		\end{center}	

\end{frame}

%%%%%%%%%%%%%%%%%%%%%%%%%%%%%%%%%%%%%%%%%%%%%%%%%%%%%%%%%%%
\begin{frame}[fragile]\frametitle{Types of Muscle Movements}

      \begin{itemize}
		\item Flexion: the joint angle gets smaller, e.g., bending the elbow.
		\item Extension: the joint angle gets larger, e.g., straightening the elbow.
		\item Abduction: a limb moves away from the midline, e.g., raising the arms sideways.
		\item Adduction: a limb moves towards the midline, e.g., bringing the legs together.
		\item Elevation: moving up, e.g., shrugging the shoulders.
		\item Depression: moving down, e.g., lowering the shoulders.
		\item Pronation: turning the palms down.
		\item Supination: turning the palms up.
		\item Rotation: turning around an axis, inwards or outwards.
		\item Sphincter action: narrowing or widening an opening.
	  \end{itemize}

  
		{\tiny (Ref: https://med.libretexts.org/Bookshelves/Anatomy\_and\_Physiology/}		
  
\end{frame}

% %%%%%%%%%%%%%%%%%%%%%%%%%%%%%%%%%%%%%%%%%%%%%%%%%%%%%%%%%%%
% \begin{frame}[fragile]\frametitle{Types of Muscle Movements}

		% \begin{center}
		% \includegraphics[width=0.4\linewidth,keepaspectratio]{muscletypes}
				
		% \end{center}	

  
		% {\tiny (Ref: https://med.libretexts.org/Bookshelves/Anatomy\_and\_Physiology/}		
  
% \end{frame}

%%%%%%%%%%%%%%%%%%%%%%%%%%%%%%%%%%%%%%%%%%%%%%%%%%%%%%%%%%%
\begin{frame}[fragile]\frametitle{Flexors}

 Muscles  which  bend  one  segment 
upon another.

		\begin{center}
		% FIG ycb1-fig-flexors | now: images/ycb1_flexors.png | final: books/YCB1/images/new/ycb1_flexors_new.pdf | register: books/YCB1/docs/image_register.csv
		% PROMPT SP-DIAGRAM (common text: books/YCB1/docs/image_system_prompt.md) | Used in block ycb1-pd-sukshma-vyayama, frame 'Flexors'. Muscle action: Flexors. Simple stick figure in side view, two overlapping positions (start and end) in light and dark line, the working muscle group drawn as a thick coloured line along the limb, one thin arrow for the movement. Label only with the single word 'Flexors'. Aspect ratio 4:3. Check before use: qualified subject reviewer.
		% NEW-PIC approved by the author (2026-11-27); the original is commented out below it
		\includegraphics[width=0.6\linewidth,keepaspectratio]{new/ycb1_flexors_new.pdf}
		% original, replaced (author approved 2026-11-27): \includegraphics[width=0.6\linewidth,keepaspectratio]{ycb1_flexors}
				
		\end{center}	

  
		% credit of the replaced original: {\tiny (Ref: YCB Official Guidebook)}
		{\tiny (Line diagram drawn for this deck.)}
  
\end{frame}


%%%%%%%%%%%%%%%%%%%%%%%%%%%%%%%%%%%%%%%%%%%%%%%%%%%%%%%%%%%
\begin{frame}[fragile]\frametitle{Extensors}

 Muscles that straighten a joint (the opposite of flexion).

		\begin{center}
		% FIG ycb1-fig-extensors | now: images/ycb1_extensors.png | final: books/YCB1/images/new/ycb1_extensors_new.pdf | register: books/YCB1/docs/image_register.csv
		% PROMPT SP-DIAGRAM (common text: books/YCB1/docs/image_system_prompt.md) | Used in block ycb1-pd-sukshma-vyayama, frame 'Extensors'. Muscle action: Extensors. Simple stick figure in side view, two overlapping positions (start and end) in light and dark line, the working muscle group drawn as a thick coloured line along the limb, one thin arrow for the movement. Label only with the single word 'Extensors'. Aspect ratio 4:3. Check before use: qualified subject reviewer.
		% NEW-PIC approved by the author (2026-11-27); the original is commented out below it
		\includegraphics[width=0.6\linewidth,keepaspectratio]{new/ycb1_extensors_new.pdf}
		% original, replaced (author approved 2026-11-27): \includegraphics[width=0.6\linewidth,keepaspectratio]{ycb1_extensors}
				
		\end{center}	

  
		% credit of the replaced original: {\tiny (Ref: YCB Official Guidebook)}
		{\tiny (Line diagram drawn for this deck.)}
  
\end{frame}

%%%%%%%%%%%%%%%%%%%%%%%%%%%%%%%%%%%%%%%%%%%%%%%%%%%%%%%%%%%
\begin{frame}[fragile]\frametitle{Abductors}

Muscles that move a part away from the midline of the body.

		\begin{center}
		% FIG ycb1-fig-abductors | now: images/ycb1_abductors.png | final: books/YCB1/images/new/ycb1_abductors_new.pdf | register: books/YCB1/docs/image_register.csv
		% PROMPT SP-DIAGRAM (common text: books/YCB1/docs/image_system_prompt.md) | Used in block ycb1-pd-sukshma-vyayama, frame 'Abductors'. Muscle action: Abductors. Simple stick figure in side view, two overlapping positions (start and end) in light and dark line, the working muscle group drawn as a thick coloured line along the limb, one thin arrow for the movement. Label only with the single word 'Abductors'. Aspect ratio 4:3. Check before use: qualified subject reviewer.
		% NEW-PIC approved by the author (2026-11-27); the original is commented out below it
		\includegraphics[width=0.6\linewidth,keepaspectratio]{new/ycb1_abductors_new.pdf}
		% original, replaced (author approved 2026-11-27): \includegraphics[width=0.6\linewidth,keepaspectratio]{ycb1_abductors}
				
		\end{center}	

  
		% credit of the replaced original: {\tiny (Ref: YCB Official Guidebook)}
		{\tiny (Line diagram drawn for this deck.)}
  
\end{frame}

%%%%%%%%%%%%%%%%%%%%%%%%%%%%%%%%%%%%%%%%%%%%%%%%%%%%%%%%%%%
\begin{frame}[fragile]\frametitle{Adductors}

 Muscles that move a part back towards the midline (the opposite of abduction).

		\begin{center}
		% FIG ycb1-fig-adductors | now: images/ycb1_adductors.png | final: books/YCB1/images/new/ycb1_adductors_new.pdf | register: books/YCB1/docs/image_register.csv
		% PROMPT SP-DIAGRAM (common text: books/YCB1/docs/image_system_prompt.md) | Used in block ycb1-pd-sukshma-vyayama, frame 'Adductors'. Muscle action: Adductors. Simple stick figure in side view, two overlapping positions (start and end) in light and dark line, the working muscle group drawn as a thick coloured line along the limb, one thin arrow for the movement. Label only with the single word 'Adductors'. Aspect ratio 4:3. Check before use: qualified subject reviewer.
		% NEW-PIC approved by the author (2026-11-27); the original is commented out below it
		\includegraphics[width=0.6\linewidth,keepaspectratio]{new/ycb1_adductors_new.pdf}
		% original, replaced (author approved 2026-11-27): \includegraphics[width=0.6\linewidth,keepaspectratio]{ycb1_adductors}
				
		\end{center}	

  
		% credit of the replaced original: {\tiny (Ref: YCB Official Guidebook)}
		{\tiny (Line diagram drawn for this deck.)}
  
\end{frame}


%%%%%%%%%%%%%%%%%%%%%%%%%%%%%%%%%%%%%%%%%%%%%%%%%%%%%%%%%%%
\begin{frame}[fragile]\frametitle{Rotators}

Muscles which rotate a part.

		\begin{center}
		% FIG ycb1-fig-rotators | now: images/ycb1_rotators.png | final: books/YCB1/images/new/ycb1_rotators_new.pdf | register: books/YCB1/docs/image_register.csv
		% PROMPT SP-DIAGRAM (common text: books/YCB1/docs/image_system_prompt.md) | Used in block ycb1-pd-sukshma-vyayama, frame 'Rotators'. Muscle action: Rotators. Simple stick figure in side view, two overlapping positions (start and end) in light and dark line, the working muscle group drawn as a thick coloured line along the limb, one thin arrow for the movement. Label only with the single word 'Rotators'. Aspect ratio 4:3. Check before use: qualified subject reviewer.
		% NEW-PIC approved by the author (2026-11-27); the original is commented out below it
		\includegraphics[width=0.6\linewidth,keepaspectratio]{new/ycb1_rotators_new.pdf}
		% original, replaced (author approved 2026-11-27): \includegraphics[width=0.6\linewidth,keepaspectratio]{ycb1_rotators}
				
		\end{center}	

  
		% credit of the replaced original: {\tiny (Ref: YCB Official Guidebook)}
		{\tiny (Line diagram drawn for this deck.)}
  
\end{frame}



%%%%%%%%%%%%%%%%%%%%%%%%%%%%%%%%%%%%%%%%%%%%%%%%%%%%%%%%%%%
\begin{frame}[fragile]\frametitle{Muscle Breathing}
\begin{columns}
    \begin{column}[T]{0.7\linewidth}
      \begin{itemize}
		\item Muscles need energy to contract.
		\item The energy comes from breaking down glucose, using oxygen.
		\item Aerobic: with oxygen; for steady, long activity.
		\item Example: dance, distance running.
		\item Anaerobic: with little oxygen; for fast, heavy effort.
		\item Anaerobic work forms lactic acid.
		\item Example: weight training, sprints.
		\item After hard effort the body takes in extra oxygen to restore its energy stores and clear the lactic acid.
		\item Oxygen debt: this extra oxygen taken in after the effort.
	  \end{itemize}
    \end{column}
    \begin{column}[T]{0.3\linewidth}
		\begin{center}
		% FIG ycb1-fig-article-respiration-muscles-wc1 | now: images/article-Respiration-Muscles-WC1.jpg | final: books/YCB1/images/new/article-Respiration-Muscles-WC1_new.pdf | register: books/YCB1/docs/image_register.csv
		% PROMPT SP-DIAGRAM (common text: books/YCB1/docs/image_system_prompt.md) | Used in block ycb1-hl-human-body-introduction, frame 'Muscle Breathing'. Anatomy diagram: Muscle Breathing. A simple schematic in front or side view, whichever shows the structure best; main parts as flat shapes with thin outlines, grayscale-safe. Labels: only the terms listed under this figure in the notes block (list them here before generating); each label short, horizontal, with a leader line. One idea only; leave out finer detail. Aspect ratio 4:3. Check before use: qualified subject reviewer (health and anatomy facts).
		% NEW-PIC (uncomment to show the new picture beside the original; comment the original only after review) \includegraphics[width=0.35\linewidth,keepaspectratio]{new/article-Respiration-Muscles-WC1_new.pdf}
		\includegraphics[width=\linewidth,keepaspectratio]{article-Respiration-Muscles-WC1.jpg}
				
		{\tiny (Ref: https://www.brainkart.com/article/Respiration-Muscles\_21121/)}
		\end{center}	
    \end{column}
  \end{columns}
\end{frame}

%%%%%%%%%%%%%%%%%%%%%%%%%%%%%%%%%%%%%%%%%%%%%%%%%%%%%%%%%%%
\begin{frame}[fragile,shrink=6]\frametitle{Cardiovascular System}
\begin{columns}
    \begin{column}[T]{0.6\linewidth}
      \begin{itemize}
		\item The system carries nutrients, gases, waste and hormones.
		\item Blood is made of:
			\begin{itemize}
				\item Plasma, the largest share (about 54\%)
				\item Red blood cells (about 45\%)
				\item White blood cells and platelets (about 1\%)
			\end{itemize}
		\item Red cells: carry oxygen and carbon dioxide.
		\item Red cells are made in the bone marrow and live about 120 days.
		\item Anaemia: haemoglobin below the WHO level for age and sex.
		\item White cells: defend the body (immunity); live about 30 hours to 25 days.
		\item Platelets: plug damaged vessels to stop bleeding.
		\item Platelets live about 8 to 10 days.
	  \end{itemize}
    \end{column}
    \begin{column}[T]{0.4\linewidth}
		\begin{center}
		% FIG ycb1-fig-cardiovascularsystem | now: images/CardiovascularSystem.jpg | final: books/YCB1/images/new/CardiovascularSystem_new.pdf | register: books/YCB1/docs/image_register.csv
		% PROMPT SP-DIAGRAM (common text: books/YCB1/docs/image_system_prompt.md) | Used in block ycb1-hl-human-body-introduction, frame 'Cardiovascular System'. Anatomy diagram: Cardiovascular System. A simple schematic in front or side view, whichever shows the structure best; main parts as flat shapes with thin outlines, grayscale-safe. Labels: only the terms listed under this figure in the notes block (list them here before generating); each label short, horizontal, with a leader line. One idea only; leave out finer detail. Aspect ratio 4:3. Check before use: qualified subject reviewer (health and anatomy facts).
		% NEW-PIC (uncomment to show the new picture beside the original; comment the original only after review) \includegraphics[width=0.35\linewidth,keepaspectratio]{new/CardiovascularSystem_new.pdf}
		\includegraphics[width=\linewidth,keepaspectratio]{CardiovascularSystem.jpg}
		
				
		{\tiny (Ref: https://www.meresearch.org.uk/taking-heart-1/)}		
		\end{center}	
    \end{column}
  \end{columns}
\end{frame}

%%%%%%%%%%%%%%%%%%%%%%%%%%%%%%%%%%%%%%%%%%%%%%%%%%%%%%%%%%%
\begin{frame}[fragile]\frametitle{Blood Vessels \& Heart}

      \begin{itemize}
		\item Blood vessels carry blood around the body.
		\item Arteries carry blood away from the heart.
		\item Arteries branch into arterioles, then capillaries, where exchange happens.
		\item Capillaries join into venules, which join into veins.
		\item Veins bring blood back to the heart.
		\item Systemic circulation: serves the body.
		\item Pulmonary circulation: serves the lungs.
		\item Heart: a muscular pump.
		\item Four chambers: right and left atria, right and left ventricles.
		\item Atria receive blood; ventricles pump it out.
		\item Valves stop backflow: tricuspid on the right, bicuspid (mitral) on the left.
	  \end{itemize}
  
\end{frame}

%%%%%%%%%%%%%%%%%%%%%%%%%%%%%%%%%%%%%%%%%%%%%%%%%%%%%%%%%%%
\begin{frame}[fragile]\frametitle{Blood Vessels \& Heart}

		\begin{center}
		% FIG ycb1-fig-heart-diagram | now: images/Heart-diagram.png | final: books/YCB1/images/new/Heart-diagram_new.pdf | register: books/YCB1/docs/image_register.csv
		% PROMPT SP-DIAGRAM (common text: books/YCB1/docs/image_system_prompt.md) | Used in block ycb1-hl-human-body-introduction, frame 'Blood Vessels and Heart'. Anatomy diagram: Blood Vessels and Heart. A simple schematic in front or side view, whichever shows the structure best; main parts as flat shapes with thin outlines, grayscale-safe. Labels: only the terms listed under this figure in the notes block (list them here before generating); each label short, horizontal, with a leader line. One idea only; leave out finer detail. Aspect ratio 4:3. Check before use: qualified subject reviewer (health and anatomy facts).
		% NEW-PIC (uncomment to show the new picture beside the original; comment the original only after review) \includegraphics[width=0.35\linewidth,keepaspectratio]{new/Heart-diagram_new.pdf}
		\includegraphics[width=0.6\linewidth,keepaspectratio]{Heart-diagram}
				
		{\tiny (Ref: https://www.meresearch.org.uk/taking-heart-1/)}
		\end{center}	

\end{frame}


%%%%%%%%%%%%%%%%%%%%%%%%%%%%%%%%%%%%%%%%%%%%%%%%%%%%%%%%%%%
\begin{frame}[fragile]\frametitle{Functions of Circulatory System}

      \begin{itemize}
		\item Transport, blood circulation.
		\item Protection, immunity.
		\item Homeostasis.
	  \end{itemize}

\end{frame}

%%%%%%%%%%%%%%%%%%%%%%%%%%%%%%%%%%%%%%%%%%%%%%%%%%%%%%%%%%%
\begin{frame}[fragile]\frametitle{Respiratory System}

      \begin{itemize}
		\item Respiration: exchange of oxygen and carbon dioxide.
		\item In the lungs: oxygen passes into the capillaries, $CO_2$ into the air sacs (alveoli).
		\item In the body: gases are exchanged in the capillaries near the cells.
		\item Respiratory tract: the path of air to and from the lungs.
		\item Nose: filters, warms and moistens air; organ of smell.
		\item Pharynx: the passage behind the mouth and nose, leading to the larynx.
		\item Larynx: the voice box.
		\item Trachea: the windpipe, held open by rings of cartilage.
		\item Bronchi and bronchioles: branching airways ending in alveoli, where gases are exchanged.
		\item Lungs: the right lung has three lobes, the left has two.
		\item Mucus lining: traps dust and irritants and moves them up to the pharynx.
		\item Sinuses: air spaces near the nose that can get blocked (sinusitis).
	  \end{itemize}

\end{frame}

%%%%%%%%%%%%%%%%%%%%%%%%%%%%%%%%%%%%%%%%%%%%%%%%%%%%%%%%%%%
\begin{frame}[fragile]\frametitle{Respiratory System}
\
		\begin{center}
		% FIG ycb1-fig-respiratory-system | now: images/respiratory-system.jpg | final: books/YCB1/images/new/respiratory-system_new.pdf | register: books/YCB1/docs/image_register.csv
		% PROMPT SP-DIAGRAM (common text: books/YCB1/docs/image_system_prompt.md) | Used in block ycb1-hl-human-body-introduction, frame 'Respiratory System'. Anatomy diagram: Respiratory System. A simple schematic in front or side view, whichever shows the structure best; main parts as flat shapes with thin outlines, grayscale-safe. Labels: only the terms listed under this figure in the notes block (list them here before generating); each label short, horizontal, with a leader line. One idea only; leave out finer detail. Aspect ratio 4:3. Check before use: qualified subject reviewer (health and anatomy facts).
		% NEW-PIC (uncomment to show the new picture beside the original; comment the original only after review) \includegraphics[width=0.35\linewidth,keepaspectratio]{new/respiratory-system_new.pdf}
		\includegraphics[width=0.45\linewidth,keepaspectratio]{respiratory-system.jpg}
		\end{center}	

\end{frame}
%%%%%%%%%%%%%%%%%%%%%%%%%%%%%%%%%%%%%%%%%%%%%%%%%%%%%%%%%%%
\begin{frame}[fragile]\frametitle{Functions of Respiratory System}

      \begin{itemize}
		\item Exchange of gases
		\item Maintaining pH balance
		\item Speech production.
	  \end{itemize}

\end{frame}
	
%%%%%%%%%%%%%%%%%%%%%%%%%%%%%%%%%%%%%%%%%%%%%%%%%%%%%%%%%%%
\begin{frame}[fragile]\frametitle{Muscles of Respiration}
\begin{columns}
    \begin{column}[T]{0.5\linewidth}
      \begin{itemize}
		\item Diaphragm: Dome-shaped muscle below lungs; separates chest and abdominal cavities.
		\item Intercostal Muscles: Located between ribs; lift rib cage for inhalation, lower it for exhalation.
		\item Accessory Muscles: Neck muscles attached to collarbone; assist in clavicular breathing.
		\item Muscles of Expiration: Abdominal muscles; used for forceful exhalation.
	  \end{itemize}
    \end{column}
    \begin{column}[T]{0.5\linewidth}
		\begin{center}
		% FIG ycb1-fig-breathing | now: images/breathing.jpg | final: books/YCB1/images/new/breathing_new.pdf | register: books/YCB1/docs/image_register.csv
		% PROMPT SP-DIAGRAM (common text: books/YCB1/docs/image_system_prompt.md) | Used in block ycb1-hl-human-body-introduction, frame 'Muscles of Respiration'. Anatomy diagram: Muscles of Respiration. A simple schematic in front or side view, whichever shows the structure best; main parts as flat shapes with thin outlines, grayscale-safe. Labels: only the terms listed under this figure in the notes block (list them here before generating); each label short, horizontal, with a leader line. One idea only; leave out finer detail. Aspect ratio 4:3. Check before use: qualified subject reviewer (health and anatomy facts).
		% NEW-PIC (uncomment to show the new picture beside the original; comment the original only after review) \includegraphics[width=0.35\linewidth,keepaspectratio]{new/breathing_new.pdf}
		\includegraphics[width=\linewidth,keepaspectratio]{breathing.jpg}
		
		{\tiny (Ref: https://step1.medbullets.com/respiratory/117007/muscles-of-respiration)}
		\end{center}	
    \end{column}
  \end{columns}
\end{frame}

%%%%%%%%%%%%%%%%%%%%%%%%%%%%%%%%%%%%%%%%%%%%%%%%%%%%%%%%%%%
\begin{frame}[fragile]\frametitle{Digestive System}

      \begin{itemize}
		\item Digestion: breaking food down into simple parts (glucose, fatty acids, amino acids).
		\item Alimentary canal: a long muscular tube with a moist lining.
		\item Food moves along by peristalsis, waves of muscle contraction.
		\item Mouth: chewing; saliva starts the digestion of carbohydrates.
		\item Oesophagus: a passage to the stomach; no digestion or absorption.
		\item Stomach: churns food and makes hydrochloric acid; protein digestion starts; no absorption.
		\item Small intestine (about 6 m): duodenum, jejunum, ileum; main digestion and absorption.
		\item Villi in the small intestine widen the surface for absorption.
		\item Large intestine: takes up water and forms faeces.
		\item Its parts: ascending, transverse and descending colon, rectum, anus.
	  \end{itemize}

\end{frame}

%%%%%%%%%%%%%%%%%%%%%%%%%%%%%%%%%%%%%%%%%%%%%%%%%%%%%%%%%%%
\begin{frame}[fragile]\frametitle{Digestive System}

		\begin{center}
		% FIG ycb1-fig-digestive-system | now: images/digestive-system.jpg | final: books/YCB1/images/new/digestive-system_new.pdf | register: books/YCB1/docs/image_register.csv
		% PROMPT SP-DIAGRAM (common text: books/YCB1/docs/image_system_prompt.md) | Used in block ycb1-hl-human-body-introduction, frame 'Digestive System'. Anatomy diagram: Digestive System. A simple schematic in front or side view, whichever shows the structure best; main parts as flat shapes with thin outlines, grayscale-safe. Labels: only the terms listed under this figure in the notes block (list them here before generating); each label short, horizontal, with a leader line. One idea only; leave out finer detail. Aspect ratio 4:3. Check before use: qualified subject reviewer (health and anatomy facts).
		% NEW-PIC (uncomment to show the new picture beside the original; comment the original only after review) \includegraphics[width=0.35\linewidth,keepaspectratio]{new/digestive-system_new.pdf}
		\includegraphics[width=0.45\linewidth,keepaspectratio]{digestive-system.jpg}
		
		{\tiny (Ref: https://my.clevelandclinic.org/health/body/7041-digestive-system)}
		\end{center}	

\end{frame}

%%%%%%%%%%%%%%%%%%%%%%%%%%%%%%%%%%%%%%%%%%%%%%%%%%%%%%%%%%%
\begin{frame}[fragile]\frametitle{Excretory System}
\begin{columns}
    \begin{column}[T]{0.5\linewidth}
      \begin{itemize}
		\item Kidneys: two bean-shaped organs that filter the blood; each has about a million filtering units (nephrons).
		\item Ureters: muscular tubes that carry urine to the bladder by peristalsis.
		\item Bladder: a hollow organ that stores urine; the urge comes at about 300 to 500 ml.
		\item Urethra: the tube that carries urine out of the body.
	  \end{itemize}
    \end{column}
    \begin{column}[T]{0.5\linewidth}
		\begin{center}
		% FIG ycb1-fig-human-excretory-system | now: images/Human-Excretory-System.jpg | final: books/YCB1/images/new/Human-Excretory-System_new.pdf | register: books/YCB1/docs/image_register.csv
		% PROMPT SP-DIAGRAM (common text: books/YCB1/docs/image_system_prompt.md) | Used in block ycb1-hl-human-body-introduction, frame 'Excretory System'. Anatomy diagram: Excretory System. A simple schematic in front or side view, whichever shows the structure best; main parts as flat shapes with thin outlines, grayscale-safe. Labels: only the terms listed under this figure in the notes block (list them here before generating); each label short, horizontal, with a leader line. One idea only; leave out finer detail. Aspect ratio 4:3. Check before use: qualified subject reviewer (health and anatomy facts).
		% NEW-PIC (uncomment to show the new picture beside the original; comment the original only after review) \includegraphics[width=0.35\linewidth,keepaspectratio]{new/Human-Excretory-System_new.pdf}
		\includegraphics[width=\linewidth,keepaspectratio]{Human-Excretory-System.jpg}
		
		{\tiny (Ref: https://www.geeksforgeeks.org/diagram-of-excretory-system/)}
		\end{center}	
    \end{column}
  \end{columns}
\end{frame}
%%%%%%%%%%%%%%%%%%%%%%%%%%%%%%%%%%%%%%%%%%%%%%%%%%%%%%%%%%%
\begin{frame}[fragile]\frametitle{Functions of Excretory System}

      \begin{itemize}
		\item Eliminate waste from the body.
		\item Regulate blood volume and blood pressure.
		\item Control levels of electrolytes and metabolites
		\item Regulate blood pH.

	  \end{itemize}

\end{frame}

%%%%%%%%%%%%%%%%%%%%%%%%%%%%%%%%%%%%%%%%%%%%%%%%%%%%%%%%%%%
\begin{frame}[fragile]\frametitle{Endocrine System}
\begin{columns}
    \begin{column}[T]{0.5\linewidth}
      \begin{itemize}
		\item The endocrine system controls body activities through hormones.
		\item Hormones: chemical messengers released straight into the blood.
		\item Each hormone acts on its own target organs.
		\item Needed only in small amounts; some, such as thyroid hormones, are stored.
		\item Kinds: water-soluble (proteins, amines) and fat-soluble (steroids).
		\item Too much or too little of a hormone disturbs health.
	  \end{itemize}
    \end{column}
    \begin{column}[T]{0.5\linewidth}
		\begin{center}
		% FIG ycb1-fig-endocrine-system | now: images/endocrine-system.jpg | final: books/YCB1/images/new/endocrine-system_new.pdf | register: books/YCB1/docs/image_register.csv
		% PROMPT SP-DIAGRAM (common text: books/YCB1/docs/image_system_prompt.md) | Used in block ycb1-hl-human-body-introduction, frame 'Endocrine System'. Anatomy diagram: Endocrine System. A simple schematic in front or side view, whichever shows the structure best; main parts as flat shapes with thin outlines, grayscale-safe. Labels: only the terms listed under this figure in the notes block (list them here before generating); each label short, horizontal, with a leader line. One idea only; leave out finer detail. Aspect ratio 4:3. Check before use: qualified subject reviewer (health and anatomy facts).
		% NEW-PIC (uncomment to show the new picture beside the original; comment the original only after review) \includegraphics[width=0.35\linewidth,keepaspectratio]{new/endocrine-system_new.pdf}
		\includegraphics[width=0.8\linewidth,keepaspectratio]{endocrine-system.jpg}
		\end{center}	
    \end{column}
  \end{columns}
\end{frame}

%%%%%%%%%%%%%%%%%%%%%%%%%%%%%%%%%%%%%%%%%%%%%%%%%%%%%%%%%%%
\begin{frame}[fragile]\frametitle{Endocrine Glands: Hypothalamus and Pituitary}

      \begin{itemize}
		\item Hypothalamus: Directs pituitary gland.
		  \begin{itemize}
		    \item Releasing Hormone (RH): Stimulates pituitary hormone release.
		    \item Inhibiting Hormone (IH): Stops pituitary hormone release.
		  \end{itemize}
		\item Pituitary Gland: Master gland; regulates other endocrine glands.
		  \begin{itemize}
		    \item Growth Hormone (GH): Promotes growth.
		    \item Follicle Stimulating Hormone (FSH): Stimulates egg and sperm formation.
		    \item Luteinizing Hormone (LH): Stimulates corpus luteum and hormone production.
		    \item Prolactin: Milk secretion.
		    \item Thyroid Stimulating Hormone (TSH): Stimulates thyroid.
		    \item Adrenocorticotropic Hormone (ACTH): Stimulates adrenal glands.
		    \item Antidiuretic Hormone (ADH): Increases water reabsorption.
		    \item Oxytocin: Uterine contractions.
		  \end{itemize}
	  \end{itemize}

\end{frame}

%%%%%%%%%%%%%%%%%%%%%%%%%%%%%%%%%%%%%%%%%%%%%%%%%%%%%%%%%%%
\begin{frame}[fragile]\frametitle{Endocrine Glands: Hypothalamus and Pituitary}

		\begin{center}
		% FIG ycb1-fig-pituitary-gland | now: images/pituitary-gland.jpg | final: books/YCB1/images/new/pituitary-gland_new.pdf | register: books/YCB1/docs/image_register.csv
		% PROMPT SP-DIAGRAM (common text: books/YCB1/docs/image_system_prompt.md) | Used in block ycb1-hl-human-body-introduction, frame 'Endocrine Glands: Hypothalamus and Pituitary'. Anatomy diagram: Endocrine Glands: Hypothalamus and Pituitary. A simple schematic in front or side view, whichever shows the structure best; main parts as flat shapes with thin outlines, grayscale-safe. Labels: only the terms listed under this figure in the notes block (list them here before generating); each label short, horizontal, with a leader line. One idea only; leave out finer detail. Aspect ratio 4:3. Check before use: qualified subject reviewer (health and anatomy facts).
		% NEW-PIC (uncomment to show the new picture beside the original; comment the original only after review) \includegraphics[width=0.35\linewidth,keepaspectratio]{new/pituitary-gland_new.pdf}
		\includegraphics[width=0.5\linewidth,keepaspectratio]{pituitary-gland.jpg}
		
		{\tiny (Ref: https://my.clevelandclinic.org/health/body/21459-pituitary-gland)}
		\end{center}	

\end{frame}


%%%%%%%%%%%%%%%%%%%%%%%%%%%%%%%%%%%%%%%%%%%%%%%%%%%%%%%%%%%
\begin{frame}[fragile]\frametitle{Endocrine Glands: Pineal, Thyroid, and Parathyroid}
\begin{columns}
    \begin{column}[T]{0.5\linewidth}
      \begin{itemize}
		\item Pineal gland: makes melatonin; sets the sleep rhythm.
		\item Thyroid: makes thyroxine and calcitonin; controls metabolism and calcium.
		\item Parathyroid glands: keep calcium in balance.
	  \end{itemize}
    \end{column}
    \begin{column}[T]{0.5\linewidth}
		\begin{center}
		% FIG ycb1-fig-endocrine | now: images/endocrine.jpg | final: books/YCB1/images/new/endocrine_new.pdf | register: books/YCB1/docs/image_register.csv
		% PROMPT SP-DIAGRAM (common text: books/YCB1/docs/image_system_prompt.md) | Used in block ycb1-hl-human-body-introduction, frame 'Endocrine Glands: Pineal, Thyroid, and Parathyroid'. Anatomy diagram: Endocrine Glands: Pineal, Thyroid, and Parathyroid. A simple schematic in front or side view, whichever shows the structure best; main parts as flat shapes with thin outlines, grayscale-safe. Labels: only the terms listed under this figure in the notes block (list them here before generating); each label short, horizontal, with a leader line. One idea only; leave out finer detail. Aspect ratio 4:3. Check before use: qualified subject reviewer (health and anatomy facts).
		% NEW-PIC (uncomment to show the new picture beside the original; comment the original only after review) \includegraphics[width=0.35\linewidth,keepaspectratio]{new/endocrine_new.pdf}
		\includegraphics[width=\linewidth,keepaspectratio]{endocrine}
		\end{center}	
    \end{column}
  \end{columns}
\end{frame}

%%%%%%%%%%%%%%%%%%%%%%%%%%%%%%%%%%%%%%%%%%%%%%%%%%%%%%%%%%%
\begin{frame}[fragile]\frametitle{Overview of Reproductive System}
\begin{columns}
    \begin{column}[T]{0.5\linewidth}
      \begin{itemize}
		\item Essential for species survival.
		\item Humans procreate via sexual reproduction.
		\item Gametes: sperm (male) and egg (female).
		\item Fertilization forms a zygote.
		\item Zygote develops into an embryo, then a fetus.
	  \end{itemize}
    \end{column}
    \begin{column}[T]{0.5\linewidth}
		\begin{center}
		% FIG ycb1-fig-reprosys | now: images/reprosys.jpg | final: books/YCB1/images/new/reprosys_new.pdf | register: books/YCB1/docs/image_register.csv
		% PROMPT SP-DIAGRAM (common text: books/YCB1/docs/image_system_prompt.md) | Used in block ycb1-hl-human-body-introduction, frame 'Overview of Reproductive System'. Anatomy diagram: Overview of Reproductive System. A simple schematic in front or side view, whichever shows the structure best; main parts as flat shapes with thin outlines, grayscale-safe. Labels: only the terms listed under this figure in the notes block (list them here before generating); each label short, horizontal, with a leader line. One idea only; leave out finer detail. Aspect ratio 4:3. Check before use: qualified subject reviewer (health and anatomy facts).
		% NEW-PIC (uncomment to show the new picture beside the original; comment the original only after review) \includegraphics[width=0.35\linewidth,keepaspectratio]{new/reprosys_new.pdf}
		\includegraphics[width=\linewidth,keepaspectratio]{reprosys}

		
		{\tiny (Ref: https://www.britannica.com/science/human-reproductive-system)}		
		\end{center}	
    \end{column}
  \end{columns}
\end{frame}

%%%%%%%%%%%%%%%%%%%%%%%%%%%%%%%%%%%%%%%%%%%%%%%%%%%%%%%%%%%
\begin{frame}[fragile]\frametitle{Male Reproductive System}
\begin{columns}
    \begin{column}[T]{0.5\linewidth}
      \begin{itemize}
		\item Testes: Oval-shaped, produce sperm.
		\item Scrotum: Sac that holds testes.
		\item Seminal Vesicles: Produce seminal fluid.
		\item Prostate Gland: Adds fluids to semen.
		\item Penis: Passes urine and semen.
	  \end{itemize}
    \end{column}
    \begin{column}[T]{0.5\linewidth}
		\begin{center}
		% FIG ycb1-fig-human-male-reproductive-system-en | now: images/Human_male_reproductive_system_en.jpg | final: books/YCB1/images/new/Human_male_reproductive_system_en_new.pdf | register: books/YCB1/docs/image_register.csv
		% PROMPT SP-DIAGRAM (common text: books/YCB1/docs/image_system_prompt.md) | Used in block ycb1-hl-human-body-introduction, frame 'Male Reproductive System'. Anatomy diagram: Male Reproductive System. A simple schematic in front or side view, whichever shows the structure best; main parts as flat shapes with thin outlines, grayscale-safe. Labels: only the terms listed under this figure in the notes block (list them here before generating); each label short, horizontal, with a leader line. One idea only; leave out finer detail. Aspect ratio 4:3. Check before use: qualified subject reviewer (health and anatomy facts).
		% NEW-PIC (uncomment to show the new picture beside the original; comment the original only after review) \includegraphics[width=0.35\linewidth,keepaspectratio]{new/Human_male_reproductive_system_en_new.pdf}
		\includegraphics[width=\linewidth,keepaspectratio]{Human_male_reproductive_system_en}
		
		{\tiny (Ref: https://simple.wikipedia.org/wiki/Male\_reproductive\_system)}
		\end{center}	
    \end{column}
  \end{columns}
\end{frame}

%%%%%%%%%%%%%%%%%%%%%%%%%%%%%%%%%%%%%%%%%%%%%%%%%%%%%%%%%%%
\begin{frame}[fragile]\frametitle{Female Reproductive System}
\begin{columns}
    \begin{column}[T]{0.5\linewidth}
      \begin{itemize}
		\item Key organs: Ovaries, oviducts, uterus, vagina.
		\item Functions: Egg production, fertilization, embryo development.
		\item Ovaries: Produce and mature eggs.
		\item Oviducts: Site of fertilization.
		\item Uterus: Houses and nurtures the embryo.
	  \end{itemize}
    \end{column}
    \begin{column}[T]{0.5\linewidth}
		\begin{center}
		% FIG ycb1-fig-female-reproductive-system | now: images/female-reproductive-system.jpg | final: books/YCB1/images/new/female-reproductive-system_new.pdf | register: books/YCB1/docs/image_register.csv
		% PROMPT SP-DIAGRAM (common text: books/YCB1/docs/image_system_prompt.md) | Used in block ycb1-hl-human-body-introduction, frame 'Female Reproductive System'. Anatomy diagram: Female Reproductive System. A simple schematic in front or side view, whichever shows the structure best; main parts as flat shapes with thin outlines, grayscale-safe. Labels: only the terms listed under this figure in the notes block (list them here before generating); each label short, horizontal, with a leader line. One idea only; leave out finer detail. Aspect ratio 4:3. Check before use: qualified subject reviewer (health and anatomy facts).
		% NEW-PIC (uncomment to show the new picture beside the original; comment the original only after review) \includegraphics[width=0.35\linewidth,keepaspectratio]{new/female-reproductive-system_new.pdf}
		\includegraphics[width=\linewidth,keepaspectratio]{female-reproductive-system.jpg}
		
		{\tiny (Ref: https://www.healthdirect.gov.au/female-reproductive-system)}		
		\end{center}	
    \end{column}
  \end{columns}
\end{frame}

%%%%%%%%%%%%%%%%%%%%%%%%%%%%%%%%%%%%%%%%%%%%%%%%%%%%%%%%%%%
\begin{frame}[fragile]\frametitle{Overview of Nervous System}
\begin{columns}
    \begin{column}[T]{0.5\linewidth}
      \begin{itemize}
		\item Coordinates and controls body actions.
		\item Neuron: basic functional unit.
		\item Consists of CNS and PNS.
		\item CNS: Brain and spinal cord.
		\item PNS: Nerves connecting CNS to body.
	  \end{itemize}
    \end{column}
    \begin{column}[T]{0.5\linewidth}
		\begin{center}
		% FIG ycb1-fig-overview-nervious-system | now: images/overview_nervious_system.jpg | final: books/YCB1/images/new/overview_nervious_system_new.pdf | register: books/YCB1/docs/image_register.csv
		% PROMPT SP-DIAGRAM (common text: books/YCB1/docs/image_system_prompt.md) | Used in block ycb1-hl-human-body-introduction, frame 'Overview of Nervous System'. Anatomy diagram: Overview of Nervous System. A simple schematic in front or side view, whichever shows the structure best; main parts as flat shapes with thin outlines, grayscale-safe. Labels: only the terms listed under this figure in the notes block (list them here before generating); each label short, horizontal, with a leader line. One idea only; leave out finer detail. Aspect ratio 4:3. Check before use: qualified subject reviewer (health and anatomy facts).
		% NEW-PIC approved by the author (2026-11-27); the original is commented out below it
		\includegraphics[width=\linewidth,keepaspectratio]{new/overview_nervious_system_new.pdf}
		% original, replaced (author approved 2026-11-27): \includegraphics[width=0.6\linewidth,keepaspectratio]{overview_nervious_system}
		
		% credit of the replaced original: {\tiny (Ref: https://neupsykey.com/overview-of-the-nervous-system/)}
		{\tiny (Line diagram drawn for this deck.)}
		
		\end{center}	
    \end{column}
  \end{columns}
\end{frame}

%%%%%%%%%%%%%%%%%%%%%%%%%%%%%%%%%%%%%%%%%%%%%%%%%%%%%%%%%%%
\begin{frame}[fragile]\frametitle{Central Nervous System: Brain}
\begin{columns}
    \begin{column}[T]{0.5\linewidth}
      \begin{itemize}
		\item Brain: Protected by skull and meninges.
		\item Cerebrum: Largest part, divided into lobes.
		\item Cerebellum: Coordinates movements and balance.
		\item Medulla Oblongata: Controls vital functions.
		\item Diencephalon: Includes hypothalamus and thalamus.
	  \end{itemize}
    \end{column}
    \begin{column}[T]{0.5\linewidth}
		\begin{center}
		% FIG ycb1-fig-basic-brain-anatomy | now: images/Basic-brain-anatomy.jpg | final: books/YCB1/images/new/Basic-brain-anatomy_new.pdf | register: books/YCB1/docs/image_register.csv
		% PROMPT SP-DIAGRAM (common text: books/YCB1/docs/image_system_prompt.md) | Used in block ycb1-hl-human-body-introduction, frame 'Central Nervous System: Brain'. Anatomy diagram: Central Nervous System: Brain. A simple schematic in front or side view, whichever shows the structure best; main parts as flat shapes with thin outlines, grayscale-safe. Labels: only the terms listed under this figure in the notes block (list them here before generating); each label short, horizontal, with a leader line. One idea only; leave out finer detail. Aspect ratio 4:3. Check before use: qualified subject reviewer (health and anatomy facts).
		% NEW-PIC (uncomment to show the new picture beside the original; comment the original only after review) \includegraphics[width=0.35\linewidth,keepaspectratio]{new/Basic-brain-anatomy_new.pdf}
		\includegraphics[width=\linewidth,keepaspectratio]{Basic-brain-anatomy.jpg}
		
		{\tiny (Ref: https://qbi.uq.edu.au/brain/brain-anatomy/central-nervous-system-brain-and-spinal-cord)}
		
		\end{center}	
    \end{column}
  \end{columns}
\end{frame}

%%%%%%%%%%%%%%%%%%%%%%%%%%%%%%%%%%%%%%%%%%%%%%%%%%%%%%%%%%%
\begin{frame}[fragile]\frametitle{Central Nervous System: Spinal Cord}
\begin{columns}
    \begin{column}[T]{0.5\linewidth}
      \begin{itemize}
		\item Runs from the medulla oblongata down to the lumbar region.
		\item Covered by meninges.
		\item Facilitates reflex actions.
		\item Conduction of sensory and motor impulses.
		\item Key role in communication between brain and body.
	  \end{itemize}
    \end{column}
    \begin{column}[T]{0.5\linewidth}
		\begin{center}
		% FIG ycb1-fig-spinal-chord | now: images/Spinal_chord.png | final: books/YCB1/images/new/Spinal_chord_new.pdf | register: books/YCB1/docs/image_register.csv
		% PROMPT SP-DIAGRAM (common text: books/YCB1/docs/image_system_prompt.md) | Used in block ycb1-hl-human-body-introduction, frame 'Central Nervous System: Spinal Cord'. Anatomy diagram: Central Nervous System: Spinal Cord. A simple schematic in front or side view, whichever shows the structure best; main parts as flat shapes with thin outlines, grayscale-safe. Labels: only the terms listed under this figure in the notes block (list them here before generating); each label short, horizontal, with a leader line. One idea only; leave out finer detail. Aspect ratio 4:3. Check before use: qualified subject reviewer (health and anatomy facts).
		% NEW-PIC (uncomment to show the new picture beside the original; comment the original only after review) \includegraphics[width=0.35\linewidth,keepaspectratio]{new/Spinal_chord_new.pdf}
		\includegraphics[width=\linewidth,keepaspectratio]{Spinal_chord}
		
		{\tiny (Ref: https://www.ck12.org/biology/central-nervous-system/lesson/central-nervous-system-ms-ls/)}
		\end{center}	
    \end{column}
  \end{columns}
\end{frame}

%%%%%%%%%%%%%%%%%%%%%%%%%%%%%%%%%%%%%%%%%%%%%%%%%%%%%%%%%%%
\begin{frame}[fragile]\frametitle{Peripheral Nervous System: Overview}
\begin{columns}
    \begin{column}[T]{0.5\linewidth}
      \begin{itemize}
		\item Includes all nerves outside CNS.
		\item Connects CNS to limbs and organs.
		\item Divided into Somatic and Autonomic systems.
		\item Somatic: Controls voluntary movements.
		\item Autonomic: Regulates involuntary functions.
	  \end{itemize}
    \end{column}
    \begin{column}[T]{0.5\linewidth}
		\begin{center}
		% FIG ycb1-fig-peripheral-nervous-system | now: images/Peripheral_Nervous_System.jpg | final: books/YCB1/images/new/Peripheral_Nervous_System_new.pdf | register: books/YCB1/docs/image_register.csv
		% PROMPT SP-DIAGRAM (common text: books/YCB1/docs/image_system_prompt.md) | Used in block ycb1-hl-human-body-introduction, frame 'Peripheral Nervous System: Overview'. Anatomy diagram: Peripheral Nervous System: Overview. A simple schematic in front or side view, whichever shows the structure best; main parts as flat shapes with thin outlines, grayscale-safe. Labels: only the terms listed under this figure in the notes block (list them here before generating); each label short, horizontal, with a leader line. One idea only; leave out finer detail. Aspect ratio 4:3. Check before use: qualified subject reviewer (health and anatomy facts).
		% NEW-PIC (uncomment to show the new picture beside the original; comment the original only after review) \includegraphics[width=0.35\linewidth,keepaspectratio]{new/Peripheral_Nervous_System_new.pdf}
		\includegraphics[width=\linewidth,keepaspectratio]{Peripheral_Nervous_System}
		
		{\tiny (Ref: https://www.simplypsychology.org/peripheral-nervous-system.html)}		
		\end{center}	
    \end{column}
  \end{columns}
\end{frame}

%%%%%%%%%%%%%%%%%%%%%%%%%%%%%%%%%%%%%%%%%%%%%%%%%%%%%%%%%%%
\begin{frame}[fragile]\frametitle{Somatic Nervous System (SNS)}
\begin{columns}
    \begin{column}[T]{0.5\linewidth}
      \begin{itemize}
		\item Sensory nerves: Carry impulses to CNS.
		\item Motor nerves: Carry impulses from CNS.
		\item 12 pairs of cranial nerves.
		\item 31 pairs of spinal nerves.
		\item Controls voluntary movements.
	  \end{itemize}
    \end{column}
    \begin{column}[T]{0.5\linewidth}
		\begin{center}
		% FIG ycb1-fig-somatic-nervous-system | now: images/Somatic_nervous_system.jpg | final: books/YCB1/images/new/Somatic_nervous_system_new.pdf | register: books/YCB1/docs/image_register.csv
		% PROMPT SP-DIAGRAM (common text: books/YCB1/docs/image_system_prompt.md) | Used in block ycb1-hl-human-body-introduction, frame 'Somatic Nervous System (SNS)'. Anatomy diagram: Somatic Nervous System (SNS). A simple schematic in front or side view, whichever shows the structure best; main parts as flat shapes with thin outlines, grayscale-safe. Labels: only the terms listed under this figure in the notes block (list them here before generating); each label short, horizontal, with a leader line. One idea only; leave out finer detail. Aspect ratio 4:3. Check before use: qualified subject reviewer (health and anatomy facts).
		% NEW-PIC (uncomment to show the new picture beside the original; comment the original only after review) \includegraphics[width=0.35\linewidth,keepaspectratio]{new/Somatic_nervous_system_new.pdf}
		\includegraphics[width=0.8\linewidth,keepaspectratio]{Somatic_nervous_system}
		
		{\tiny (Ref: https://www.simplypsychology.org/somatic-nervous-system.html)}
		\end{center}	
    \end{column}
  \end{columns}
\end{frame}

%%%%%%%%%%%%%%%%%%%%%%%%%%%%%%%%%%%%%%%%%%%%%%%%%%%%%%%%%%%
\begin{frame}[fragile]\frametitle{Autonomic Nervous System (ANS)}
\begin{columns}
    \begin{column}[T]{0.5\linewidth}
      \begin{itemize}
		\item Regulates involuntary actions.
		\item Sympathetic: 'Fight or flight' response.
		\item Parasympathetic: 'Rest and digest' response.
		\item Controls internal organs.
		\item Includes sympathetic and parasympathetic chains.
	  \end{itemize}
    \end{column}
    \begin{column}[T]{0.5\linewidth}
		\begin{center}
		% FIG ycb1-fig-autonomic-nervous-system | now: images/Autonomic_Nervous_System.jpg | final: books/YCB1/images/new/Autonomic_Nervous_System_new.pdf | register: books/YCB1/docs/image_register.csv
		% PROMPT SP-DIAGRAM (common text: books/YCB1/docs/image_system_prompt.md) | Used in block ycb1-hl-human-body-introduction, frame 'Autonomic Nervous System (ANS)'. Anatomy diagram: Autonomic Nervous System (ANS). A simple schematic in front or side view, whichever shows the structure best; main parts as flat shapes with thin outlines, grayscale-safe. Labels: only the terms listed under this figure in the notes block (list them here before generating); each label short, horizontal, with a leader line. One idea only; leave out finer detail. Aspect ratio 4:3. Check before use: qualified subject reviewer (health and anatomy facts).
		% NEW-PIC approved by the author (2026-11-27); the original is commented out below it
		\includegraphics[width=\linewidth,keepaspectratio]{new/Autonomic_Nervous_System_new.pdf}
		% original, replaced (author approved 2026-11-27): \includegraphics[width=\linewidth,keepaspectratio]{Autonomic_Nervous_System}
		
		% credit of the replaced original: {\tiny (Ref: https://www.simplypsychology.org/autonomic-nervous-system.html)}
		{\tiny (Line diagram drawn for this deck.)}
		\end{center}	
    \end{column}
  \end{columns}
\end{frame}


%%%%%%%%%%%%%%%%%%%%%%%%%%%%%%%%%%%%%%%%%%%%%%%%%%%%%%%%%%%
% Dropped 2026-11-27: the screenshot tables credit asanas with cures (food poisoning, diabetes, grey hair), against the no-cure rule.
% \begin{frame}[fragile]\frametitle{Asana Benefits}
% 
% 
% 		\begin{center}
% 		% FIG ycb1-fig-asanabenefits-1 | now: images/ycb1_asanabenefits_1.png | final: no image file (see register) | register: books/YCB1/docs/image_register.csv
% 		% PROMPT SP-NONE (common text: books/YCB1/docs/image_system_prompt.md) | Do not generate a picture. This is a screenshot of a table. Retype it as a LaTeX table (deck) and a Markdown table (notes) from the checked notes block ycb1-pd-yogasana-practice, in own words where it is a list of facts. Used in block ycb1-pd-yogasana-practice, frame 'Asana Benefits'.
% 		\includegraphics[width=\linewidth,keepaspectratio]{ycb1_asanabenefits_1}
% 				
% 		\end{center}	
% 
%   
% 		{\tiny (Ref: YCB Official Guidebook)}		
%   
% \end{frame}



%%%%%%%%%%%%%%%%%%%%%%%%%%%%%%%%%%%%%%%%%%%%%%%%%%%%%%%%%%%
% Dropped 2026-11-27: the screenshot tables credit asanas with cures (food poisoning, diabetes, grey hair), against the no-cure rule.
% \begin{frame}[fragile]\frametitle{Asana Benefits}
% 
% 
% 		\begin{center}
% 		% FIG ycb1-fig-asanabenefits-2 | now: images/ycb1_asanabenefits_2.png | final: no image file (see register) | register: books/YCB1/docs/image_register.csv
% 		% PROMPT SP-NONE (common text: books/YCB1/docs/image_system_prompt.md) | Do not generate a picture. This is a screenshot of a table. Retype it as a LaTeX table (deck) and a Markdown table (notes) from the checked notes block ycb1-pd-yogasana-practice, in own words where it is a list of facts. Used in block ycb1-pd-yogasana-practice, frame 'Asana Benefits'.
% 		\includegraphics[width=\linewidth,keepaspectratio]{ycb1_asanabenefits_2}
% 		
% 		% FIG ycb1-fig-asanabenefits-3 | now: images/ycb1_asanabenefits_3.png | final: no image file (see register) | register: books/YCB1/docs/image_register.csv
% 		% PROMPT SP-NONE (common text: books/YCB1/docs/image_system_prompt.md) | Do not generate a picture. This is a screenshot of a table. Retype it as a LaTeX table (deck) and a Markdown table (notes) from the checked notes block ycb1-pd-yogasana-practice, in own words where it is a list of facts. Used in block ycb1-pd-yogasana-practice, frame 'Asana Benefits'.
% 		\includegraphics[width=\linewidth,keepaspectratio]{ycb1_asanabenefits_3}
% 		
% 				
% 		\end{center}	
% 
%   
% 		{\tiny (Ref: YCB Official Guidebook)}		
%   
% \end{frame}


%%%%%%%%%%%%%%%%%%%%%%%%%%%%%%%%%%%%%%%%%%%%%%%%%%%%%%%%%%%%%%%%%%%%%%%%%%%%%%%%%%
\begin{frame}[fragile]\frametitle{}
\begin{center}
{\Large 3.2  Meaning and Means of health promotion and role of Yoga  in health promotion}
\end{center}
\end{frame}

%%%%%%%%%%%%%%%%%%%%%%%%%%%%%%%%%%%%%%%%%%%%%%%%%%%%%%%%%%%%%%%%%%%%%%%%%%%%%%%%%%
\begin{frame}[fragile]\frametitle{Health Promotion}
    \textbf{Definition of Health:}
    \begin{itemize}
        \item According to the World Health Organization (WHO), health is defined as \textit{“a state of complete physical, mental, and social well-being and not merely the absence of disease or infirmity.”}
    \end{itemize}

    \textbf{Health Promotion:}
    \begin{itemize}
        \item Health promotion involves \textit{encouraging healthy lifestyles} and educating people to increase \textit{awareness} and \textit{well-being}.
        \item Objective: To promote \textit{healthy living} and \textit{prevent disease}.
    \end{itemize}
\end{frame}

%%%%%%%%%%%%%%%%%%%%%%%%%%%%%%%%%%%%%%%%%%%%%%%%%%%%%%%%%%%%%%%%%%%%%%%%%%%%%%%%%%
\begin{frame}[fragile,shrink=6]\frametitle{Role of Yoga in Health Promotion}
    \textbf{Yoga and Health Benefits:}
    \begin{itemize}
        \item Areas of health that yoga touches:
        \begin{itemize}
            % Removed: the notes leave out heart and circulation claims.
            % \item \textbf{Cardiovascular Fitness} - Regular practice builds general fitness.
            \item \textbf{Strength and balance:} regular practice builds muscle strength and balance.
            \item \textbf{Joint mobility:} the joints move more freely.
            % Removed: stamina and endurance are not in the notes.
            % \item \textbf{Stamina and Body Balance} - Builds endurance and stability.
        \end{itemize}
        \item \textbf{Shatkriyas (शट्क्रियास):} internal cleansing practices, learnt from a qualified teacher.
        \item \textbf{Pranayama (प्राणायाम):} works on the breath and on \textit{prana} (प्राण).
        \item \textbf{Dhyana (ध्यान), meditation:} builds concentration.
    \end{itemize}

    \textbf{Experiential Learning:}
    \begin{itemize}
        \item Yoga is learnt by doing. Steady practice gives a better grasp of one's own well-being.
        \item Try different paths, such as \textbf{Bhakti Yoga (भक्ति योग)}, \textbf{Karma Yoga (कर्म योग)} and \textbf{Ishvara Pranidhana (ईश्वर प्रणिधान)}, and see what suits you.
        \item Keep up steady practice, as Patanjali advises (\textit{abhyasa, अभ्यास}), and watch its effects.
    \end{itemize}
\end{frame}

%%%%%%%%%%%%%%%%%%%%%%%%%%%%%%%%%%%%%%%%%%%%%%%%%%%%%%%%%%%
\begin{frame}[fragile]\frametitle{Effects of Hatha Yoga Practices}
      \begin{itemize}
        \item Enhances flexibility in tendons, muscles, and spine.
        \item Builds strength, balance and joint mobility.
        \item Stiffness and pain can lessen with regular practice.
        \item Brings attention to breath and body.
        \item Supports better posture.
        \item Eases tension in joints and muscles.
        \item Loosens tight areas like neck and shoulders.
        \item Promotes mental well-being alongside physical health.
      \end{itemize}
\end{frame}

%%%%%%%%%%%%%%%%%%%%%%%%%%%%%%%%%%%%%%%%%%%%%%%%%%%%%%%%%%%
\begin{frame}[fragile]\frametitle{Limitations and Contraindications of Yoga Practices}
      \begin{itemize}
        \item Learn the contraindications before starting practice.
        \item Yoga is mainly for prevention and well-being.
        \item People differ: not every practice suits everyone, so do not compare yourself with others.
        % Removed: a treatment claim; the notes keep to prevention and well-being.
        % \item More effective for functional disorders than for organic conditions.
        % Removed: a treatment claim about a named disease; left out of the notes.
        % \item Not a cure for conditions like cancer; helps in managing symptoms and improving strength.
        \item Its effects need support from a sound diet and lifestyle.
        \item It works alongside medical care, never instead of it.
      \end{itemize}
\end{frame}


%%%%%%%%%%%%%%%%%%%%%%%%%%%%%%%%%%%%%%%%%%%%%%%%%%%%%%%%%%%%%%%%%%%%%%%%%%%%%%%%%%
\begin{frame}[fragile]\frametitle{}
\begin{center}
{\Large 3.3 Yogic positive Attitudes}
\end{center}
\end{frame}

%%%%%%%%%%%%%%%%%%%%%%%%%%%%%%%%%%%%%%%%%%%%%%%%%%%%%%%%%%%
\begin{frame}[fragile]\frametitle{Yogic Positive Attitudes in Yogasutra: Chitta Prasadana}

मैत्रीकरुणामुदितोपेक्षाणां सुखदुःखपुण्यापुण्यविषयाणां भावनातश्चित्तप्रसादनम् ॥ १.३३ ॥

      \begin{itemize}
        \item \textbf{Maitri:} friendliness towards the happy.
        \item \textbf{Karuna:} compassion towards those who suffer.
        \item \textbf{Mudita:} gladness towards the virtuous, those who do well.
        \item \textbf{Upeksha:} even-minded acceptance towards those who do wrong, neither joining in nor condemning.
        \item \textbf{Clarity of Mind:} the sutra names these attitudes as the way to a clear, calm mind (\textit{chitta prasadana}).
        \item \textbf{Inner Peace:} peace, calm and serenity of mind.
        \item \textbf{Emotional Balance:} steadier emotional responses.
        \item \textbf{Self-Improvement:} personal growth may follow (interpretation).
        \item \textbf{Harmonious Relationships:} kinder attitudes may ease dealings with others (interpretation).
        \item \textbf{Mindful Awareness:} attention to one's own attitudes in daily life (interpretation).
      \end{itemize}

\end{frame}

%%%%%%%%%%%%%%%%%%%%%%%%%%%%%%%%%%%%%%%%%%%%%%%%%%%%%%%%%%%%%%%%%%%%%%%%%%%%%%%%%%
\begin{frame}[fragile]\frametitle{}
\begin{center}
{\Large 3.4 Concept of Bhavas}
\end{center}
\end{frame}

% %%%%%%%%%%%%%%%%%%%%%%%%%%%%%%%%%%%%%%%%%%%%%%%%%%%%%%%%%%%
% \begin{frame}[fragile]\frametitle{Concept of Bhavas}

      % \begin{itemize}
        % \item Bhavas are essential attitudes for mental and emotional well-being.
        % \item They include Dharma, Jnana, Vairagya, and Aishwarya.
        % \item Each Bhava contributes uniquely to balanced living and personal growth.
      % \end{itemize}

% \end{frame}

% %%%%%%%%%%%%%%%%%%%%%%%%%%%%%%%%%%%%%%%%%%%%%%%%%%%%%%%%%%%
% \begin{frame}[fragile]\frametitle{Dharma (Duty)}

      % \begin{itemize}
        % \item \textbf{Sanskrit Verse:} 
        
        % योगिश्चत्तवृित्तिनरोधः   ॥   १ . २॥ \\
        % Yogachittavrittinirodhah || 1.2
        
        % \item Dharma is the highest duty to maintain a balanced mind.
        % \item It involves using Yoga techniques like Asana, Pranayama, and Meditation.
        % \item Aim to keep a positive state even amidst life’s ups and downs.

        % \item \textbf{Sanskrit Verse:} 
        
        % दृग्दशर्नशक्त्योरेकात्मतेवािस्मता   ॥   २ . ६॥ \\
        % Drk-Daesanasaktyoh-Ekatmate-Iva-Asmita || 2.6
        
        % \item Gain understanding and knowledge before taking action.
        % \item Learn from self-experience; view learning as personal growth.
        % \item Avoid ego by recognizing that knowledge is personal and subjective.
      % \end{itemize}	  

% \end{frame}

% %%%%%%%%%%%%%%%%%%%%%%%%%%%%%%%%%%%%%%%%%%%%%%%%%%%%%%%%%%%
% \begin{frame}[fragile]\frametitle{Vairagya (Objectivity) and Aishwarya (Self-Reliance)}

      % \begin{itemize}
        % \item \textbf{Vairagya (Objectivity):}
        
        % तत्परं   पुरुषख्यातेगुर्णवैतृष्ण्यम्   ॥   १ . १६॥ \\
        % Tatparam Purusakhyateh Gunavaitrshnyam || 1.16
        
        % \item Cultivate detachment and a witness-like attitude.
        % \item Surrender ego and remain objective in life.
        % \item \textbf{Aishwarya (Self-Reliance):}
        
        % तत्र   स्थितौ   यत्नोऽभ्यासः   ॥   १ . १३॥ \\
        % Tatra sthitau-yatnah abhyasah || 1.13
        
        % \item Achieve self-reliance and perfection through consistent practice.
        % \item Increases willpower, confidence, and joy.
      % \end{itemize}

% \end{frame}

%%%%%%%%%%%%%%%%%%%%%%%%%%%%%%%%%%%%%%%%%%%%%%%%%%%%%%%%%%%%%%%%%%%%%%%%%%%%%%%%%%
\begin{frame}[fragile]\frametitle{Concept of Bhava}
    \textbf{Bhava (भाव)} can be understood as:
    \begin{itemize}
        \item A certain \textbf{state of mind}
        \item A \textbf{feeling}, \textbf{emotion}, or \textbf{attitude}
    \end{itemize}

    \textbf{Yogic Science:}
    \begin{itemize}
        \item Yogic practices work on the mind by bringing up a particular \textit{bhava} within oneself.
        \item Workshop teaching, not a sourced fact: some postures are said to bring up certain bhavas.
    \end{itemize}

    \textbf{Types of Bhava:}
    \begin{itemize}
        \item \textbf{Positive Bhavas:} Dharma Bhava (धर्म भाव), Jnana Bhava (ज्ञान भाव), Vairagya Bhava (वैराग्य भाव), and Aishvarya Bhava (ऐश्वर्य भाव)
        \item \textbf{Negative Bhavas:} Adharma Bhava (अधर्म भाव), Raga Bhava (राग भाव), Dvesha Bhava (द्वेष भाव), and Ajnana Bhava (अज्ञान भाव)
    \end{itemize}
\end{frame}


%%%%%%%%%%%%%%%%%%%%%%%%%%%%%%%%%%%%%%%%%%%%%%%%%%%%%%%%%%%%%%%%%%%%%%%%%%%%%%%%%%
\begin{frame}[fragile,shrink=6]\frametitle{Understanding and Practicing Bhava}
    \textbf{Negative Bhavas:}
    \begin{itemize}
        \item \textbf{Adharma Bhava (अधर्म भाव)} - Unrighteousness
        \item \textbf{Raga Bhava (राग भाव)} - Attachment
        \item \textbf{Dvesha Bhava (द्वेष भाव)} - Aversion
        \item \textbf{Ajnana Bhava (अज्ञान भाव)} - Ignorance
    \end{itemize}
    
    \textbf{Positive Bhavas:}
    \begin{itemize}
        \item \textbf{Dharma Bhava (धर्म भाव)} - Discipline and Duty
        \item \textbf{Jnana Bhava (ज्ञान भाव)} - Knowledge and Clarity
        \item \textbf{Vairagya Bhava (वैराग्य भाव)} - Detachment and Letting Go
        \item \textbf{Aishvarya Bhava (ऐश्वर्य भाव)} - Inner strength and power
    \end{itemize}

    \textbf{Practices to Enhance Positive Bhavas (workshop teaching, no named source):}
    \begin{itemize}
        \item \textbf{Dharma Bhava:} Padmasana (पद्मासन), Sukhasana (सुखासन), and the Yamas (यम) and Niyamas (नियम).
        \item \textbf{Jnana Bhava:} balancing postures, Kapalabhati (कपालभाति) and Trataka (त्राटक).
        \item \textbf{Vairagya Bhava:} forward bending postures.
        \item \textbf{Aishvarya Bhava:} building self-awareness and introspection.
    \end{itemize}
\end{frame}


%%%%%%%%%%%%%%%%%%%%%%%%%%%%%%%%%%%%%%%%%%%%%%%%%%%%%%%%%%%%%%%%%%%%%%%%%%%%%%%%%%
\begin{frame}[fragile]\frametitle{}
\begin{center}
{\Large 3.5 Dincharya and Ritucharya with respect to Yogic lifestyle}
\end{center}
\end{frame}

% %%%%%%%%%%%%%%%%%%%%%%%%%%%%%%%%%%%%%%%%%%%%%%%%%%%%%%%%%%%%%%%%%%%%%%%%%%%%%%%%%%
% \begin{frame}[fragile]\frametitle{Ayurvedic Principles: Dinacharya and Ritucharya}

      % \begin{itemize}
        % \item \textbf{Relative Health Concepts:} What is healthy or unhealthy depends on individual factors like age, constitution, and health conditions.
        % \item \textbf{Context Matters:} Diet and lifestyle recommendations vary based on factors such as season, disease state, and geographical location.
        % \item \textbf{Dinacharya (Daily Regimen):} Ayurveda prescribes daily routines to balance Doshas and prevent disorders.
        % \item \textbf{Ritucharya (Seasonal Regimen):} Seasonal guidelines help align lifestyle with environmental changes to maintain health.
        % \item \textbf{Preventive Focus:} Ayurveda emphasizes prevention over cure; following Dinacharya and Ritucharya enhances immunity and life quality.
        % \item \textbf{Long-Term Benefits:} Adhering to these regimens reduces risk of future disorders and improves overall life expectancy.
        % \item \textbf{Scientific Basis:} Ancient Ayurvedic principles, tested over centuries, remain relevant and beneficial globally.
        % \item \textbf{Universal Relevance:} Basic Ayurvedic principles apply across different climates and regions, adapting tools but maintaining core concepts.
      % \end{itemize}

% \end{frame}

% %%%%%%%%%%%%%%%%%%%%%%%%%%%%%%%%%%%%%%%%%%%%%%%%%%%%%%%%%%%%%%%%%%%%%%%%%%%%%%%%%%
% \begin{frame}[fragile]\frametitle{Ritucharya: Seasonal Regimen}

      % \begin{itemize}
        % \item \textbf{Similarity-Dissimilarity Doctrine:} 
        
        % “A similar material enriches the similar in the body, while a dissimilar material depletes its counterpart.”
        
        % \item \textbf{Seasonal Impact on Doshas:} 
        % \begin{itemize}
            % \item Cold weather increases Vata Dosha and decreases Pitta Dosha.
            % \item Dry conditions can aggravate Kapha Dosha.
        % \end{itemize}
        % \item \textbf{Balancing External Influences:} 
        % \begin{itemize}
            % \item Ritucharya helps balance the internal environment against seasonal changes.
            % \item Lifestyle and diet are adapted to minimize adverse effects.
        % \end{itemize}
        % \item \textbf{Dosha Reactions:} 
        % \begin{itemize}
            % \item \textbf{Sanchaya:} Accumulation of Dosha.
            % \item \textbf{Prakopa:} Aggravation of Dosha.
            % \item \textbf{Prashama:} Pacification of Dosha.
        % \end{itemize}
        % \item \textbf{Maintaining Dosha Equilibrium:} 
        % \begin{itemize}
            % \item Aim to achieve Dosha-Samya (Dosha equilibrium).
            % \item Focus on transitional periods (Ritusandhi) between seasons to prevent diseases.
        % \end{itemize}
        % \item \textbf{Health Benefits:} 
        % \begin{itemize}
            % \item Enhances immunity, prevents seasonal diseases, and improves quality of life.
        % \end{itemize}
      % \end{itemize}

% \end{frame}

% %%%%%%%%%%%%%%%%%%%%%%%%%%%%%%%%%%%%%%%%%%%%%%%%%%%%%%%%%%%%%%%%%%%%%%%%%%%%%%%%%%
% \begin{frame}[fragile]\frametitle{Ritusandhi: Transitional Period Between Seasons}

      % \begin{itemize}
        % \item \textbf{Definition:} 
        
        % Ritusandhi refers to the transitional period between two seasons, lasting about two weeks.
        
        % \item \textbf{Duration:} 
        % \begin{itemize}
            % \item First week: Last seven days of the current season.
            % \item Second week: First seven days of the upcoming season.
        % \end{itemize}
        % \item \textbf{Gradual Changes:} 
        % \begin{itemize}
            % \item Seasonal transitions are gradual, not abrupt.
            % \item Adjustments in diet and routine should be made gradually.
        % \end{itemize}
        % \item \textbf{Caution Required:} 
        % \begin{itemize}
            % \item This period is delicate; careful management of diet and routine is crucial.
            % \item Avoid practices that disrupt Dosha balance.
        % \end{itemize}
        % \item \textbf{Health Risks:} 
        % \begin{itemize}
            % \item Common health issues during Ritusandhi include colds, flu, and fever.
            % \item Hospitals often see increased admissions during this time.
        % \end{itemize}
        % \item \textbf{Preventive Measures:} 
        % \begin{itemize}
            % \item Adhere to Ayurvedic recommendations to maintain health and prevent illness.
        % \end{itemize}
      % \end{itemize}

% \end{frame}

% %%%%%%%%%%%%%%%%%%%%%%%%%%%%%%%%%%%%%%%%%%%%%%%%%%%%%%%%%%%%%%%%%%%%%%%%%%%%%%%%%%
% \begin{frame}[fragile]\frametitle{Classification of Seasons}

      % \begin{itemize}
        % \item \textbf{Samwatsara (Year):} 
        
        % A year consists of two Ayanas and six Ritus.
        
        % \item \textbf{Ayanas (Semesters):} 
        % \begin{itemize}
            % \item \textbf{Uttarayana (Northern Solstice):}
            
            % The Sun moves northward, taking moisture from the earth (Adanakala).
            
            % \item \textbf{Dakshinayana (Southern Solstice):}
            
            % The Sun moves southward, giving moisture to the earth (Visargakala).
            
        % \end{itemize}
        % \item \textbf{Uttarayana Seasons:}
        % \begin{itemize}
            % \item \textbf{Shishira:} Late winter.
            % \item \textbf{Vasanta:} Spring.
            % \item \textbf{Grishma:} Summer.
        % \end{itemize}
        % \item \textbf{Dakshinayana Seasons:}
        % \begin{itemize}
            % \item \textbf{Varsha:} Rainy season.
            % \item \textbf{Sharad:} Autumn.
            % \item \textbf{Hemanta:} Early winter.
        % \end{itemize}
        % \item \textbf{Seasonal Attributes:} 
        % \begin{itemize}
            % \item Understanding seasonal attributes helps in selecting suitable food and activities.
        % \end{itemize}
      % \end{itemize}

% \end{frame}

% %%%%%%%%%%%%%%%%%%%%%%%%%%%%%%%%%%%%%%%%%%%%%%%%%%%%%%%%%%%%%%%%%%%%%%%%%%%%%%%%%%
% \begin{frame}[fragile]\frametitle{Dinacharya: Daily Regimen}

      % \begin{itemize}
        % \item \textbf{Concept:} 
        
        % A day reflects seasonal and yearly attributes, guiding daily health practices.
        
        % \item \textbf{Purpose:} 
        % \begin{itemize}
            % \item Promotes long-term health and well-being.
        % \end{itemize}
        % \item \textbf{Daily Routine Includes:}
        % \begin{itemize}
            % \item \textbf{Morning:} Wake up, elimination, oral care, self-massage.
            % \item \textbf{Day:} Exercise, bathing, diet, social etiquette.
            % \item \textbf{Evening:} Relaxation, bedtime routine, sexual health.
        % \end{itemize}
        % \item \textbf{Guidance:} 
        % \begin{itemize}
            % \item Detailed Ayurveda instructions on practices and materials.
        % \end{itemize}
      % \end{itemize}

% \end{frame}

% %%%%%%%%%%%%%%%%%%%%%%%%%%%%%%%%%%%%%%%%%%%%%%%%%%%%%%%%%%%%%%%%%%%%%%%%%%%%%%%%%%
% \begin{frame}[fragile]\frametitle{Ayurvedic Dinacharya: Daily Cycles and Practices}

      % \begin{itemize}
        % \item \textbf{Daily Cycles:}
        % \begin{itemize}
            % \item \textbf{Sun Cycle (6:00 AM - 6:00 PM):}
            % \begin{itemize}
                % \item 6:00 AM - 10:00 AM: Kapha
                % \item 10:00 AM - 2:00 PM: Pitta
                % \item 2:00 PM - 6:00 PM: Vata
            % \end{itemize}
            % \item \textbf{Moon Cycle (6:00 PM - 6:00 AM):}
            % \begin{itemize}
                % \item 6:00 PM - 10:00 PM: Kapha
                % \item 10:00 PM - 2:00 AM: Pitta
                % \item 2:00 AM - 6:00 AM: Vata
            % \end{itemize}
        % \end{itemize}
        % \item \textbf{Recommended Practices:}
        % \begin{itemize}
            % \item Wake up, elimination, face wash, brushing teeth, tongue cleaning.
            % \item Gargle, nasal cleaning, massage, exercise, herbal scrub, bathing.
            % \item Worship, clothing \& shoes, work, sleep \& sex.
        % \end{itemize}
        % \item \textbf{Benefits:} 
        % \begin{itemize}
            % \item Connects with nature, prevents diseases, reduces stress.
            % \item Enhances digestion, promotes discipline, peace, happiness, and longevity.
        % \end{itemize}
      % \end{itemize}

% \end{frame}


%%%%%%%%%%%%%%%%%%%%%%%%%%%%%%%%%%%%%%%%%%%%%%%%%%%%%%%%%%%%%%%%%%%%%%%%%%%%%%%%%%
\begin{frame}[fragile]\frametitle{Dinacharya and Ritucharya}
    \textbf{Dinacharya (दिनचर्या)}: Daily Routine
    \begin{itemize}
        \item A fixed routine for each day.
        \item In Ayurveda, it keeps the three doshas in balance.
        \item The aim is health and the prevention of illness.
    \end{itemize}

    \textbf{Ritucharya (ऋतुचर्या )}: Seasonal Routine
    \begin{itemize}
        \item Diet and activity changed to suit the season.
        \item It protects the body from the season's effects.
        \item At the junction of two seasons, change gradually.
    \end{itemize}

    \textbf{दोष  Doshas and Daily Phases (a common modern teaching scheme):}
    \begin{itemize}
        \item \textbf{Vata (वात)}: 2:00 AM - 6:00 AM, 2:00 PM - 6:00 PM
        \item \textbf{Kapha (कफ)}: 6:00 AM - 10:00 AM, 6:00 PM - 10:00 PM
        \item \textbf{Pitta (पित्त)}: 10:00 AM - 2:00 PM, 10:00 PM - 2:00 AM
    \end{itemize}
\end{frame}

%%%%%%%%%%%%%%%%%%%%%%%%%%%%%%%%%%%%%%%%%%%%%%%%%%%%%%%%%%%%%%%%%%%%%%%%%%%%%%%%%%
\begin{frame}[fragile]\frametitle{Daily Guidelines (Dinacharya)}
    \textbf{Suggested Routine:}
    \begin{itemize}
        \item \textbf{Brahmamuhurta (ब्राह्ममुहूर्त)}: wake about one to one and a half hours before sunrise.
        \item \textbf{Morning Routine:}
        \begin{itemize}
            % Removed: this term is not in the notes' checked list of morning items.
            % \item Mouthwash (कंठधावन)
            \item Brushing the teeth (दन्तधावन)
            \item Tongue cleaning (जिव्हा-धावन, jihwa nirlekhana)
            \item Oil pulling (गण्डूष, gandusha or kavala)
            \item Washing mouth and face (मुखप्रक्षालन)
        \end{itemize}
        \item \textbf{Daily Activities:}
        \begin{itemize}
            \item Physical exercise (व्यायाम)
            \item Body massage (अभ्यंग)
            \item Bathing (स्नान)
            \item Dressing and grooming (वस्त्रादि)
        \end{itemize}
    \end{itemize}
\end{frame}

%%%%%%%%%%%%%%%%%%%%%%%%%%%%%%%%%%%%%%%%%%%%%%%%%%%%%%%%%%%%%%%%%%%%%%%%%%%%%%%%%%
\begin{frame}[fragile]\frametitle{Seasonal Guidelines (Ritucharya)}
    \textbf{Six Seasons in Two Halves:}
    \begin{itemize}
        \item \textbf{Adana Kala (आदान काल)}: the taking-away half; strength tends to fall
        \begin{itemize}
            \item Shishira (शिशिर) - Winter
            \item Vasanta (वसंत) - Spring
            \item Grishma (ग्रीष्म) - Summer
        \end{itemize}
        \item \textbf{Visarga Kala (विसर्ग काल)}: the giving half
        \begin{itemize}
            \item Varsha (वर्षा) - Monsoon
            \item Sharad (शरद्) - Autumn
            \item Hemanta (हेमन्त) - Late autumn or early winter
        \end{itemize}
    \end{itemize}

    \textbf{Guidelines:}
    \begin{itemize}
        \item Adjust food and activity to the season.
        \item For a personal routine, ask a qualified Ayurvedic practitioner.
    \end{itemize}
\end{frame}

%%%%%%%%%%%%%%%%%%%%%%%%%%%%%%%%%%%%%%%%%%%%%%%%%%%%%%%%%%%%%%%%%%%%%%%%%%%%%%%%%%
\begin{frame}[fragile]\frametitle{}
\begin{center}
{\Large 3.6 Holistic approach of Yoga towards health and Diseases}
\end{center}
\end{frame}

% %%%%%%%%%%%%%%%%%%%%%%%%%%%%%%%%%%%%%%%%%%%%%%%%%%%%%%%%%%%%%%%%%%%%%%%%%%%%%%%%%%
% \begin{frame}[fragile]\frametitle{Yoga, Health, and Awareness}

      % \begin{itemize}
        % \item \textbf{Mindset and Health:} 
        
        % "Man today is sick because he thinks he is sick. Sickness and disease have no place in a person who does not accept self-laments." 
        
        % \item \textbf{Swami Satyananda Saraswati's View:}
        % \begin{itemize}
            % \item Disease is not fate but a call for change and growth.
            % \item Yoga helps break the illusion of disease as inescapable.
        % \end{itemize}
        % \item \textbf{Yoga's Role:}
        % \begin{itemize}
            % \item Yoga and medicine complement each other in restoring health.
            % \item Yoga views disease as a teacher and guide to balance.
        % \end{itemize}
        % \item \textbf{Disease as a Teacher:}
        % \begin{itemize}
            % \item Indicates lifestyle or mental errors.
            % \item Signals the need for lifestyle changes for better health and joy.
        % \end{itemize}
      % \end{itemize}

% \end{frame}

% %%%%%%%%%%%%%%%%%%%%%%%%%%%%%%%%%%%%%%%%%%%%%%%%%%%%%%%%%%%%%%%%%%%%%%%%%%%%%%%%%%
% \begin{frame}[fragile]\frametitle{Holistic Health and the Yogic Approach}

      % \begin{itemize}
        % \item \textbf{Regaining Awareness:}
        % \begin{itemize}
            % \item Disease forces awareness of natural law transgressions.
            % \item Yogic practices restore balance and understanding.
        % \end{itemize}
        % \item \textbf{WHO Definition of Health:}
        
        % Health is physical, mental, intellectual, and spiritual well-being, not just a disease-free body.
        
        % \item \textbf{Holistic Treatment Approach:}
        % \begin{itemize}
            % \item Traditional Yoga, therapeutic Yoga, and Ayurveda address root causes of diseases.
            % \item A holistic, empathetic approach is essential in yoga practice.
        % \end{itemize}
        % \item \textbf{Yoga Teacher's Role:}
        % \begin{itemize}
            % \item Provide complete solutions with empathy for physical and mental complaints.
        % \end{itemize}
      % \end{itemize}

% \end{frame}


%%%%%%%%%%%%%%%%%%%%%%%%%%%%%%%%%%%%%%%%%%%%%%%%%%%%%%%%%%%%%%%%%%%%%%%%%%%%%%%%%%
\begin{frame}[fragile]\frametitle{Holistic Approach of Yoga Towards Health}
    \textbf{Yoga and Health:}
    \begin{itemize}
        \item Yoga is practised widely to keep overall health.
        \item Yoga aims to act before illness appears; it supports medical care and does not replace it.
        \item \textbf{Patanjali's Insight:} \textit{हेयं दुःखमनागतम्} (Heyam duhkham anagatam, Yoga Sutra 2.16): the pain not yet come is to be avoided.
    \end{itemize}

    \textbf{Prevention Aspect:}
    \begin{itemize}
        \item Yoga practices help prevent illness and keep overall health.
        \item Yoga is basically a preventive discipline; regular practice supports well-being.
    \end{itemize}
\end{frame}


%%%%%%%%%%%%%%%%%%%%%%%%%%%%%%%%%%%%%%%%%%%%%%%%%%%%%%%%%%%%%%%%%%%%%%%%%%%%%%%%%%
\begin{frame}[fragile]\frametitle{Yoga, Care and Well-being}
    \textbf{Yoga as Therapy:}
    \begin{itemize}
        \item Yoga is mainly preventive; in chronic illness consult a doctor and a yoga therapist.
        \item Breath practices are taught for relaxation and well-being, not as a treatment.
        % Removed: a note to the writer, not content for the learner.
        % \item Claims of medical benefit from practices like \textit{कपालभाति} (Kapalbhati) need a named study; none is cited here.
    \end{itemize}

    \textbf{Holistic Benefits:}
    \begin{itemize}
        \item Yoga looks after physical, mental, social and spiritual well-being.
        \item Practices aim to support overall health and ease stress, towards a balanced life.
    \end{itemize}
\end{frame}

%%%%%%%%%%%%%%%%%%%%%%%%%%%%%%%%%%%%%%%%%%%%%%%%%%%%%%%%%%%%%%%%%%%%%%%%%%%%%%%%%%
\begin{frame}[fragile]\frametitle{}
\begin{center}
{\Large 3.7 Introduction to first aid and Cardio Pulmonary Resuscitation (CPR)}
\end{center}
\end{frame}

%%%%%%%%%%%%%%%%%%%%%%%%%%%%%%%%%%%%%%%%%%%%%%%%%%%%%%%%%%%%%%%%%%%%%%%%%%%%%%%%%%
\begin{frame}[fragile]\frametitle{Introduction to First Aid}
    \textbf{What is First Aid?}
    \begin{itemize}
        \item The first help given to an injured or suddenly ill person, before an ambulance, a medical person or a healthcare facility can take over.
        \item Aims: preserve life, stop the condition getting worse, help recovery, and get the person safely to the nearest healthcare facility.
        \item Examples: applying bandages and cleaning minor cuts and scratches.
    \end{itemize}
\end{frame}

%%%%%%%%%%%%%%%%%%%%%%%%%%%%%%%%%%%%%%%%%%%%%%%%%%%%%%%%%%%%%%%%%%%%%%%%%%%%%%%%%%
\begin{frame}[fragile]\frametitle{Four Steps in Any Emergency}
    \begin{enumerate}
        \item \textbf{Make the area safe}: your own safety first.
        \item \textbf{Evaluate the person}: conscious? breathing normally?
        \item \textbf{Seek help}: send a bystander to call 112, and to come back and confirm.
        \item \textbf{Give first aid.}
    \end{enumerate}
    \textbf{When calling, have ready:} where you are, the number you call from, what has happened, and how many are hurt and how.

    Memory aid: safe, check, call, help.
\end{frame}

%%%%%%%%%%%%%%%%%%%%%%%%%%%%%%%%%%%%%%%%%%%%%%%%%%%%%%%%%%%%%%%%%%%%%%%%%%%%%%%%%%
\begin{frame}[fragile]\frametitle{Moving a Casualty, Food and Drink}
    \begin{itemize}
        \item As a rule, do not move an injured person.
        \item Move only if staying is more dangerous, the place cannot be made safe, help will not come soon, and you are not in danger yourself.
        \item Give nothing to eat or drink to a person who is badly hurt, feels sick, is getting sleepy or is losing consciousness.
        \item The manual lists a few exceptions (for example low blood sugar), only when the person is conscious.
    \end{itemize}
\end{frame}


%%%%%%%%%%%%%%%%%%%%%%%%%%%%%%%%%%%%%%%%%%%%%%%%%%%%%%%%%%%%%%%%%%%%%%%%%%%%%%%%%%
\begin{frame}[fragile,shrink=6]\frametitle{First Aid Kit and Basic Procedures}

      \begin{itemize}
        \item \textbf{Definition:} 
        
        Immediate care given to an injured or ill person until trained help takes over.

        \item \textbf{First Aid Kit Essentials (workshop list, not yet checked):}
        \begin{itemize}
            \item Adhesive bandages, gauze pads, adhesive tape.
            \item Scissors, tweezers, a thermometer.
            % Removed: antiseptic wipes, pain relievers and burn cream are not in the notes; burn cream goes against the burn advice.
            % \item Adhesive bandages, gauze pads, antiseptic wipes.
            % \item Pain relievers, burn cream, digital thermometer.
        \end{itemize}
        \item \textbf{Basic First Aid Procedures (not yet checked against the manual):}
        \begin{itemize}
            \item \textbf{Wounds:} press with a clean cloth to stop bleeding; rinse a small cut with clean water; cover with a clean dressing.
            \item \textbf{Burns:} cool under cool running water for at least 20 minutes; no ice, butter or cream.
            \item Then cover the burn loosely with a clean non-fluffy cloth or cling film.
            \item \textbf{Fractures:} keep the part still and get medical help.
        \end{itemize}
        \item \textbf{Importance of Training:}

        A quick, correct response in an emergency can save a life.
        
      \end{itemize}

\end{frame}

%%%%%%%%%%%%%%%%%%%%%%%%%%%%%%%%%%%%%%%%%%%%%%%%%%%%%%%%%%%%%%%%%%%%%%%%%%%%%%%%%%
\begin{frame}[fragile]\frametitle{Introduction to CPR}
    \textbf{What is CPR?}
    \begin{itemize}
        \item Cardiopulmonary Resuscitation (CPR) is a life-saving procedure for a person who is not breathing, or not breathing normally.
        % Removed: a survival figure that is not in the notes.
        % \item CPR can double or triple the chances of survival after cardiac arrest.
        \item \textbf{Types of CPR:}
        \begin{itemize}
            \item \textbf{Trained helper:} compressions with breaths, 30 compressions to 2 breaths.
            \item \textbf{Untrained or little trained helper:} compressions only.
            \item \textbf{Baby under one year:} breaths are always given, with the baby technique.
        \end{itemize}
    \end{itemize}
\end{frame}

%%%%%%%%%%%%%%%%%%%%%%%%%%%%%%%%%%%%%%%%%%%%%%%%%%%%%%%%%%%%%%%%%%%%%%%%%%%%%%%%%%
\begin{frame}[fragile]\frametitle{Importance of CPR Training}
    \textbf{CPR Certification:}
    \begin{itemize}
        \item In India, CPR certification is not known to be compulsory for yoga teachers (still to be checked).
        % Removed: not in the notes.
        % \item In Western countries, it might be required by some studios.
        \item CPR training is useful for everyone; it is worth taking a course.
    \end{itemize}

    \textbf{Action Items:}
    \begin{itemize}
        % Removed: no video is named in the notes.
        % \item Watch the reference video provided by the Global Association of Indian Medical Students.
        \item Learn CPR in a hands-on course of the Indian Red Cross or St John Ambulance.
    \end{itemize}
\end{frame}


%%%%%%%%%%%%%%%%%%%%%%%%%%%%%%%%%%%%%%%%%%%%%%%%%%%%%%%%%%%%%%%%%%%%%%%%%%%%%%%%%%
\begin{frame}[fragile,shrink=6]\frametitle{CPR Steps, Depth and Rate}

      \begin{itemize}
        \item \textbf{Definition:} 
        
        A life-saving technique for a person who is not breathing, or not breathing normally.

        \item \textbf{CPR Steps:}
        \begin{itemize}
            \item \textbf{Check Response:} is the person conscious?
            \item \textbf{Call for Help:} send a bystander to call 112, India's single emergency number.
            \item \textbf{Check Breathing:} tilt the head back, lift the chin; look, listen and feel for no more than 10 seconds.
            \item Do not spend time feeling for a pulse.
            \item \textbf{Chest Compression:} push in the centre of the chest, in sets of 30.
            \item \textbf{Rescue Breaths:} if trained, give 2 breaths after every 30 compressions.
        \end{itemize}
        \item \textbf{Compression Depth and Rate:}
        \begin{itemize}
            \item Depth: 5 to 6 cm in an adult.
            \item Rate: about 100 a minute, not more than 120.
        \end{itemize}
        \item \textbf{When to Perform CPR:}

        When the person is not breathing or not breathing normally. Occasional noisy gasps just after a collapse are not normal breathing.

      \end{itemize}

\end{frame}

%%%%%%%%%%%%%%%%%%%%%%%%%%%%%%%%%%%%%%%%%%%%%%%%%%%%%%%%%%%%%%%%%%%%%%%%%%%%%%%%%%
\begin{frame}[fragile]\frametitle{CPR: Doing It Well}
    \begin{itemize}
        \item Person on their back on a hard surface; kneel beside the upper arm.
        \item Heel of one hand in the centre of the chest, other hand on top, fingers locked; shoulders over the chest, arms straight.
        \item Release fully after each push, without leaning, so the chest rises.
        \item Rescue breaths: tilt the head, lift the chin, pinch the nose, seal the mouth, blow for one second.
        \item Watch the chest rise. Give 2 breaths, at most 2 attempts, then 30 compressions again.
        \item With more than one rescuer, swap every two minutes with the smallest possible pause.
    \end{itemize}
\end{frame}

%%%%%%%%%%%%%%%%%%%%%%%%%%%%%%%%%%%%%%%%%%%%%%%%%%%%%%%%%%%%%%%%%%%%%%%%%%%%%%%%%%
\begin{frame}[fragile]\frametitle{Stopping CPR and the Instructor's Role}
    \begin{itemize}
        \item Stop only when: the person shows signs of life such as breathing, a trained person or medical staff takes over, you are too exhausted, or the place becomes unsafe.
        \item If breathing returns, place the person in the recovery position; wash your hands.
        \item The instructor calls for help early, gives first aid or CPR only as trained, and hands over to medical staff.
        \item This is an introduction for the exam; learn CPR in a hands-on course (Indian Red Cross or St John Ambulance).
    \end{itemize}
\end{frame}

%%%%%%%%%%%%%%%%%%%%%%%%%%%%%%%%%%%%%%%%%%%%%%%%%%%%%%%%%%%%%%%%%%%%%%%%%%%%%%%%%%
\begin{frame}[fragile]\frametitle{Quick check}
    \begin{enumerate}
        \item What is the emergency number in India?
        \item What compression rate does current guidance give for adult CPR?
        \item What is the usual ratio of compressions to breaths for a trained helper?
    \end{enumerate}
    {\small Answers: 112; 100 to 120 a minute; 30 compressions to 2 breaths.}
\end{frame}

%%%%%%%%%%%%%%%%%%%%%%%%%%%%%%%%%%%%%%%%%%%%%%%%%%%%%%%%%%%%%%%%%%%%%%%%%%%%%%%%%%
\begin{frame}[fragile]\frametitle{}
\begin{center}
{\Large 3.8 Yogic Management of stress and its consequences}
\end{center}
\end{frame}

%%%%%%%%%%%%%%%%%%%%%%%%%%%%%%%%%%%%%%%%%%%%%%%%%%%%%%%%%%%%%%%%%%%%%%%%%%%%%%%%%%
\begin{frame}[fragile]\frametitle{Human Psyche: Modern and Yogic Concepts}

      \begin{itemize}
        \item \textbf{Psychology:}
        
        The scientific study of mental processes and behaviour; it touches family life, education and health.

        \item \textbf{Behavior:}
        \begin{itemize}
            \item \textbf{Overt Behavior:} what can be seen from outside: actions and reactions to a stimulus.
            \item \textbf{Covert Behavior:} the inner mental activity.
        \end{itemize}
        \item \textbf{Consciousness:}
        
        In the yogic view, the awareness behind mind and experience, which can expand towards super-consciousness.
        
        \item \textbf{Consciousness Expansion:}
        \begin{itemize}
            \item In the yogic view, the final goal is the union of the individual Self (Atman) with the Supreme Self (Paramatman).
        \end{itemize}
      \end{itemize}

\end{frame}

%%%%%%%%%%%%%%%%%%%%%%%%%%%%%%%%%%%%%%%%%%%%%%%%%%%%%%%%%%%%%%%%%%%%%%%%%%%%%%%%%%
\begin{frame}[fragile]\frametitle{Indian Model of Personality}

      \begin{itemize}
        \item \textbf{Upanishadic Personality Model:}
        
        The person is described as five sheaths, or \textit{Koshas} (\textit{कोश}).

        \item \textbf{Annamaya Kosha:}

        The food sheath, fed by \textit{Anna} (\textit{अन्न}, food).

        \item \textbf{Pranamaya Kosha:}

        The vital sheath, fed by \textit{Prana} (\textit{प्राण}, life energy).

        \item \textbf{Manomaya Kosha:}

        The mental sheath, fed by \textit{Pratyahara} (\textit{प्रत्याहार}, drawing in the senses).

        \item \textbf{Vijnanamaya Kosha:}

        The intellect sheath, fed by \textit{Dhyana} (\textit{ध्यान}, meditation).

        \item \textbf{Anandamaya Kosha:}

        The bliss sheath, fed by \textit{Samadhi} (\textit{समाधि}).
        
      \end{itemize}

\end{frame}


%%%%%%%%%%%%%%%%%%%%%%%%%%%%%%%%%%%%%%%%%%%%%%%%%%%%%%%%%%%%%%%%%%%%%%%%%%%%%%%%%%
\begin{frame}[fragile]\frametitle{Development of Consciousness: The Three Gunas}

      \begin{itemize}
        \item \textbf{Three Gunas:}
        
        Fundamental qualities influencing consciousness and behavior.
        
        \item \textbf{Sattva} (\textit{सत्त्व}): Stability
        \begin{itemize}
            \item Attributes: Love, compassion, honesty, and calm.
        \end{itemize}
        \item \textbf{Rajas} (\textit{रजस्}): Activation
        \begin{itemize}
            \item Attributes: Action, ambition, desire, and leadership.
        \end{itemize}
        \item \textbf{Tamas} (\textit{तमस्}): Inertia
        \begin{itemize}
            \item Attributes: Laziness, sleep, dullness, and aversion.
        \end{itemize}
        \item \textbf{Mental Functions:}
        
        \textit{Vrittis} (\textit{वृत्तिः}) and \textit{Pravritti} (\textit{प्रवृत्ति}) are manifestations of the Three Gunas.
        
      \end{itemize}

\end{frame}

%%%%%%%%%%%%%%%%%%%%%%%%%%%%%%%%%%%%%%%%%%%%%%%%%%%%%%%%%%%%%%%%%%%%%%%%%%%%%%%%%%
\begin{frame}[fragile]\frametitle{Causes of Frustrations and Psychosomatic Disorders}

      \begin{itemize}
        \item \textbf{Mind as a Conglomeration of Thoughts:}
        
        Thoughts are compared to waves on the sea.

        \item \textbf{Process of Mental Activity:}
        \begin{itemize}
            \item The senses (Indriyas) bring in information.
            \item The intellect, helped by memory, processes it.
            \item Emotion, positive or negative, then colours it.
        \end{itemize}
        \item \textbf{Negative Emotions:}

        Anger, fear, hatred and jealousy feed stress and the mind-born disorders called Adhi.

        \item \textbf{Positive Emotions:}

        Peace, contentment and happiness refresh and build up.
        
      \end{itemize}

\end{frame}

%%%%%%%%%%%%%%%%%%%%%%%%%%%%%%%%%%%%%%%%%%%%%%%%%%%%%%%%%%%%%%%%%%%%%%%%%%%%%%%%%%
\begin{frame}[fragile]\frametitle{Mental Hygiene and Its Objectives}

      \begin{itemize}
        \item \textbf{Definition:}
        
        Mental hygiene is the practice of maintaining mental health by being aware of and managing one's thoughts and emotions.
        
        \item \textbf{Objectives of Mental Hygiene:}
        \begin{itemize}
            \item Realize one's potential.
            \item Develop self-respect and respect for others.
            \item Understand and tolerate limitations of self and others.
            \item Promote harmony and happiness.
            \item Make effective adjustments in life.
            \item Know one's true self.
        \end{itemize}
      \end{itemize}

\end{frame}

%%%%%%%%%%%%%%%%%%%%%%%%%%%%%%%%%%%%%%%%%%%%%%%%%%%%%%%%%%%%%%%%%%%%%%%%%%%%%%%%%%
\begin{frame}[fragile]\frametitle{Yogic Attitudes for Mental Hygiene}

      \begin{itemize}
        \item \textbf{Pratipaksha Bhavana} (\textit{प्रतिपक्ष भावना }):
        
        Cultivating the opposite thought to cancel a disturbing one.

        \item \textbf{Anitya Bhavana} (\textit{अनित्य भावना }):

        Seeing that bodily experiences do not last, which grows detachment (\textit{वैराग्य}, Vairagya).

        \item \textbf{Sakshi Bhavana} (\textit{साक्षी भावना }):

        Watching one's own actions as a witness, which grows self-awareness and evenness of mind.
        
      \end{itemize}

\end{frame}

%%%%%%%%%%%%%%%%%%%%%%%%%%%%%%%%%%%%%%%%%%%%%%%%%%%%%%%%%%%%%%%%%%%%%%%%%%%%%%%%%%
\begin{frame}[fragile]\frametitle{Yogic Perception of Mental Health}

      \begin{itemize}
        \item \textbf{Definition:}
        
        A state of well-being in which a person knows their abilities, copes with the usual stresses, works productively and serves the community.

        \item \textbf{Patanjali's View:}
        \begin{itemize}
            \item Yoga is the stilling of the mental modifications (\textit{वृत्ति}, Vritti).
            \item The mind is stilled by \textbf{Abhyasa} (\textit{अभ्यास}, practice) and \textbf{Vairagya} (\textit{वैराग्य}, detachment).
            \item \textbf{Abhyasa} (\textit{अभ्यास}): steady, repeated practice.
            \item \textbf{Key Practices:} 
                \begin{itemize}
                    \item Pratyahara (\textit{प्रत्याहार} - Withdrawal of senses)
                    \item Dharana (\textit{धारणा} - Concentration)
                    \item Dhyana (\textit{ध्यान} - Meditation)
                    \item Samadhi (\textit{समाधि} - Self-realization)
                \end{itemize}
        \end{itemize}
      \end{itemize}

\end{frame}

%%%%%%%%%%%%%%%%%%%%%%%%%%%%%%%%%%%%%%%%%%%%%%%%%%%%%%%%%%%%%%%%%%%%%%%%%%%%%%%%%%
\begin{frame}[fragile]\frametitle{Role of Prayer and Meditation in Mental Health}

      \begin{itemize}
        \item \textbf{Prayer (the guidebook's explanations, not proven effects):}
        \begin{itemize}
            \item Described as the most widely practised form of healing.
            \item Effects offered:
                \begin{itemize}
                    \item Brings a relaxation response.
                    \item Eases the strain of trying to control everything.
                    \item Placebo response and secondary control are offered as explanations.
                    \item Fits the person's beliefs.
                    \item Raises positive feelings.
                    \item Joins body, mind and spirit.
                \end{itemize}
        \end{itemize}
        \item \textbf{Meditation Benefits:}
        \begin{itemize}
            \item \textbf{OM Meditation:} makes the mind one-pointed.
            \item Helps with focus on the task in hand.
            \item Clears information overload.
            \item A tool for self-realisation.
        \end{itemize}
        \item \textbf{Psychosocial Environment:}
        \begin{itemize}
            \item The culture and climate of a workplace.
            \item Psychosocial stress comes from dealing with other people.
        \end{itemize}
      \end{itemize}

\end{frame}

%%%%%%%%%%%%%%%%%%%%%%%%%%%%%%%%%%%%%%%%%%%%%%%%%%%%%%%%%%%%%%%%%%%%%%%%%%%%%%%%%%
\begin{frame}[fragile]\frametitle{Concept of Stress: Modern Science and Yoga}

    \begin{itemize}
        \item \textbf{Definition:}
        
        The response of a person when the demands are greater than the resources to meet them.

        \item \textbf{Types of Stress:}
        \begin{itemize}
            \item \textbf{Eustress:} mild, useful stress (e.g., excitement).
            \item \textbf{Distress:} harmful, ongoing stress (physical or psychological).
        \end{itemize}
        \item \textbf{Stress Reactions:}
        \begin{itemize}
            \item The body is prepared for ``fight or flight'': more energy, faster heart, higher blood pressure.
            \item Blood goes to the muscles and the senses sharpen.
            \item Stress that is not resolved can end in psychosomatic disorders.
        \end{itemize}
        \item \textbf{Yoga Perspective:}
        \begin{itemize}
            \item Stress is an imbalance of body and mind; its roots are the \textbf{Kleshas}.
            \item Stressors: too many obligations, negative thinking, poor handling of conflict.
        \end{itemize}
    \end{itemize}

\end{frame}

%%%%%%%%%%%%%%%%%%%%%%%%%%%%%%%%%%%%%%%%%%%%%%%%%%%%%%%%%%%%%%%%%%%%%%%%%%%%%%%%%%
\begin{frame}[fragile]\frametitle{Yogic View on Stress Management}

    \begin{itemize}
        \item \textbf{Likes and Dislikes:} 
        Strong liking and disliking upset the mind and give rise to stress (\textit{अधि}, Adhis).
        \item \textbf{Yogic Means (five terms in the guidebook):}
        \begin{itemize}
            \item \textbf{Ahara} (\textit{आहार}): right food.
            \item \textbf{Vihara} (\textit{विहार}): proper recreation and rest.
            \item \textbf{Achara} (\textit{आचार}): conduct and attitudes.
            \item \textbf{Vichara} (\textit{विचार}): positive thinking.
            \item \textbf{Vyavahara} (\textit{व्यवहार}): correct conduct and action.
        \end{itemize}
        \item \textbf{Practices:}
        Cyclic meditation (postures alternating with rest lying on the back) is reported in studies to help with stress.
        \item \textbf{Research:}
        \begin{itemize}
            \item Studies report better attention and emotional competence.
            \item Studies report lower anxiety.
        \end{itemize}
        \item \textbf{Life Management:}
        \begin{itemize}
            \item Follow \textbf{Karma Yoga} (\textit{कर्म योग}): action without attachment to its fruit.
            \item Aim: a steady mind and, in the yogic view, self-realisation.
        \end{itemize}
    \end{itemize}

\end{frame}


%%%%%%%%%%%%%%%%%%%%%%%%%%%%%%%%%%%%%%%%%%%%%%%%%%%%%%%%%%%%%%%%%%%%%%%%%%%%%%%%%%
\begin{frame}[fragile]\frametitle{}
\begin{center}
{\Large 3.9 Yoga in prevention of metabolic and respiratory disorders}
\end{center}
\end{frame}

%%%%%%%%%%%%%%%%%%%%%%%%%%%%%%%%%%%%%%%%%%%%%%%%%%%%%%%%%%%%%%%%%%%%%%%%%%%%%%%%%%
\begin{frame}[fragile,shrink=8]\frametitle{Respiratory \& Metabolic Health: Yogic Prevention}

    \begin{itemize}
        \item \textbf{Respiratory System:}
        \begin{itemize}
            \item The nose, throat, lungs, diaphragm and the muscles that help breathing.
            \item Upper and lower airways; their disorders are often linked.
        \end{itemize}
        \item \textbf{Yogic Approach:}
        \begin{itemize}
            \item \textbf{Mucus Clearing:} seen as helpful; jala neti uses warm salted water to keep the nose clear.
            \item \textbf{Cautions:} cleansing practices are learnt only from a qualified teacher.
            \item Avoid kapalabhati and other forceful breathing in heart disease, high blood pressure, epilepsy or hernia.
        \end{itemize}
        \item \textbf{Metabolic Disorders:}
        \begin{itemize}
            \item Digestive health matters for overall well-being, of body and mind.
            % Removed: a note to the writer; replaced by the shared risk factors from the notes.
            % \item Eating well and regularly is a base of a healthy lifestyle (named source needed).
            \item Shared risks: tobacco, inactivity, unhealthy diet, harmful alcohol, air pollution (WHO).
        \end{itemize}
        \item \textbf{Yogic Prevention:}
        \begin{itemize}
            \item \textbf{Healthy Eating:} right food, amount and time support well-being.
            \item \textbf{Seek Advice:} with any condition, practise only after the doctor's advice, under a qualified teacher.
            \item Yoga supports prevention; it does not replace medical treatment.
        \end{itemize}
    \end{itemize}

\end{frame}

%%%%%%%%%%%%%%%%%%%%%%%%%%%%%%%%%%%%%%%%%%%%%%%%%%%%%%%%%%%%%%%%%%%%%%%%%%%%%%%%%%
\begin{frame}[fragile,shrink=6]\frametitle{Role of Digestive Power and Yogic Management}

    \begin{itemize}
        \item \textbf{Optimal Health:}
        \begin{itemize}
            \item Needs good eating habits: the right food, amount and time.
            \item Eating for emotional comfort or from greed upsets digestion.
        \end{itemize}
        \item \textbf{Manipura Chakra} (\textit{मणिपूर चक्र}):
        \begin{itemize}
            \item At the navel region (solar plexus); stands for the digestive organs.
            \item Stands for the inner digestive fire.
        \end{itemize}
        \item \textbf{Digestive Process:}
        \begin{itemize}
            \item The fire element (\textit{अग्नि}) digests; air, water and earth support it.
        \end{itemize}
        \item \textbf{Hatha Yoga} (\textit{हठ योग}):
        \begin{itemize}
            \item Keeps the abdomen healthy with asanas (\textit{आसन}), pranayamas (\textit{प्राणायाम}) and shatkarmas (\textit{षटकर्म}).
            \item Dhauti (\textit{धौति}), nauli (\textit{नौली}) and basti (\textit{बस्ति}) are the cleansing actions for the digestive tract.
            \item The Hatha text gives them for those with excess fat or phlegm. Learn them only from a qualified teacher.
        \end{itemize}
        \item \textbf{Yogic Benefits:}
        \begin{itemize}
            \item Transforms digestion into a source of higher awareness and vitality.
        \end{itemize}
    \end{itemize}

\end{frame}


%%%%%%%%%%%%%%%%%%%%%%%%%%%%%%%%%%%%%%%%%%%%%%%%%%%%%%%%%%%%%%%%%%%%%%%%%%%%%%%%%%
\begin{frame}[fragile]\frametitle{}
\begin{center}
{\Large 3.10 Yoga for personality development}
\end{center}
\end{frame}

%%%%%%%%%%%%%%%%%%%%%%%%%%%%%%%%%%%%%%%%%%%%%%%%%%%%%%%%%%%%%%%%%%%%%%%%%%%%%%%%%%
\begin{frame}[fragile]\frametitle{Yoga for Personality Development}

    \begin{itemize}
        \item \textbf{Expected gains of regular practice (not guaranteed):}
        \item \textbf{Self-Awareness:} Enhances understanding of oneself.
        \item \textbf{Emotional Control:} Manages stress and emotions.
        \item \textbf{Discipline:} Improves focus and self-discipline.
        \item \textbf{Confidence:} Builds self-esteem through practice.
        \item \textbf{Positive Attitude:} Encourages optimism.
        \item \textbf{Relationships:} Enhances empathy and communication.
        \item \textbf{Resilience:} Strengthens mental adaptability.
        \item \textbf{Holistic Growth:} Supports overall development.
    \end{itemize}

\end{frame}


